\documentclass[12pt,a4paper]{article}

% =========================================================
% Русский язык и шрифты
% =========================================================

\usepackage[utf8]{inputenc}
\usepackage[T2A]{fontenc}
\usepackage[russian]{babel}
\usepackage{XCharter}
\usepackage{microtype}

% =========================================================
% Математика и графика
% =========================================================

\usepackage{amsmath,amssymb,amsfonts,amsthm,mathtools,mathrsfs}
\usepackage{esvect}
\usepackage{tikz}

\usetikzlibrary{
    hobby,
    arrows.meta,
    calc,
    intersections,
    decorations.pathreplacing,
    patterns.meta
}

\DeclareMathOperator{\supp}{supp}
\DeclareMathOperator{\dist}{dist}

% =========================================================
% Компактная книжная вёрстка
% =========================================================

\usepackage{geometry}
\geometry{
    top=0.8cm,
    bottom=1.0cm,
    left=1.0cm,
    right=1.0cm,
    includefoot
}

\usepackage{setspace}
\setstretch{1.06}

\setlength{\parindent}{1.25em}
\setlength{\parskip}{0pt}

\setlength{\abovedisplayskip}{6pt plus 2pt minus 2pt}
\setlength{\belowdisplayskip}{6pt plus 2pt minus 2pt}
\setlength{\abovedisplayshortskip}{3pt plus 2pt}
\setlength{\belowdisplayshortskip}{4pt plus 2pt minus 2pt}

\frenchspacing
\allowdisplaybreaks
\emergencystretch=2em

\usepackage{enumitem}

\setlist[itemize]{
    nosep,
    leftmargin=1.5em
}

\setlist[enumerate]{
    nosep,
    leftmargin=1.8em
}

% =========================================================
% Заголовки
% =========================================================

\usepackage{titlesec}

\titleformat{\section}
    {\large\bfseries}
    {\thesection.}
    {0.6em}
    {}

\titlespacing*{\section}
    {0pt}
    {1.1em}
    {0.45em}

\titleformat{\subsection}
    {\normalsize\bfseries}
    {\thesubsection.}
    {0.6em}
    {}

\titlespacing*{\subsection}
    {0pt}
    {0.85em}
    {0.25em}

\titleformat{\subsubsection}
    {\normalsize\bfseries}
    {\thesubsubsection.}
    {0.5em}
    {}

\titlespacing*{\subsubsection}
    {0pt}
    {0.65em}
    {0.2em}

% В оглавлении показываем темы и подпункты лекций.
\setcounter{secnumdepth}{3}
\setcounter{tocdepth}{3}

% =========================================================
% Колонтитулы
% =========================================================

\usepackage{fancyhdr}

\pagestyle{fancy}
\fancyhf{}
\fancyfoot[C]{\thepage}

\renewcommand{\headrulewidth}{0pt}
\renewcommand{\footrulewidth}{0pt}

\fancypagestyle{plain}{
    \fancyhf{}
    \fancyfoot[C]{\thepage}
    \renewcommand{\headrulewidth}{0pt}
    \renewcommand{\footrulewidth}{0pt}
}

% =========================================================
% Цвета
% =========================================================

\usepackage{xcolor}

\definecolor{defbg}{HTML}{F0F8FF}
\definecolor{defframe}{HTML}{5DADE2}

\definecolor{thmbg}{HTML}{F4FCF4}
\definecolor{thmframe}{HTML}{7DCEA0}

\definecolor{lembg}{HTML}{FFF9F0}
\definecolor{lemframe}{HTML}{F5B041}

\definecolor{exbg}{HTML}{F8F9F9}
\definecolor{exframe}{HTML}{BFC9CA}

\definecolor{commentbg}{HTML}{F8F5FF}
\definecolor{commentframe}{HTML}{A569BD}

\definecolor{exercisebg}{HTML}{FFF8F8}
\definecolor{exerciseframe}{HTML}{EC7063}

\definecolor{darkblue}{HTML}{154360}

% =========================================================
% Теоремные блоки
% =========================================================

\usepackage{tcolorbox}

\tcbuselibrary{
    skins,
    breakable,
    theorems
}

\tcbset{
    modernbox/.style={
        enhanced,
        breakable,
        frame hidden,
        boxrule=0pt,
        arc=1.5mm,
        top=2mm,
        bottom=2mm,
        left=3mm,
        right=3mm,
        before skip=7pt,
        after skip=7pt,
        fonttitle=\bfseries\sffamily,
        coltitle=black!80,
        attach title to upper={\par\smallskip}
    }
}

\newtcbtheorem{Definition}{Определение}{
    modernbox,
    colback=defbg,
    borderline west={3pt}{0pt}{defframe},
    coltitle=defframe!60!black
}{def}

\newtcbtheorem{Theorem}{Теорема}{
    modernbox,
    colback=thmbg,
    borderline west={3pt}{0pt}{thmframe},
    coltitle=thmframe!60!black
}{thm}

\newtcbtheorem{Lemma}{Лемма}{
    modernbox,
    colback=lembg,
    borderline west={3pt}{0pt}{lemframe},
    coltitle=lemframe!60!black
}{lem}

\newtcolorbox{Example}[1][]{
    modernbox,
    colback=exbg,
    borderline west={3pt}{0pt}{exframe},
    title={Пример #1},
    coltitle=exframe!60!black
}

\newtcolorbox{Comment}[1][]{
    modernbox,
    colback=commentbg,
    borderline west={3pt}{0pt}{commentframe},
    title={Комментарий #1},
    coltitle=commentframe!70!black
}

\newtcolorbox{Exercise}[1][]{
    modernbox,
    colback=exercisebg,
    borderline west={3pt}{0pt}{exerciseframe},
    title={Упражнение #1},
    coltitle=exerciseframe!70!black
}

\newenvironment{Proof}[1][Доказательство]
{
    \par\smallskip
    \noindent
    \textit{#1.}\enspace
    \ignorespaces
}
{
    \hfill$\blacksquare$
    \par\smallskip
}

% =========================================================
% Гиперссылки
% =========================================================

\usepackage{hyperref}

\hypersetup{
    colorlinks=true,
    linkcolor=darkblue,
    filecolor=magenta,
    urlcolor=darkblue,
    pdftitle={Квантовая механика}
}

% =========================================================
% Начало документа
% =========================================================


\begin{document}

\tableofcontents
\clearpage

% =========================================================
\section{Лекции 1--2}
% =========================================================

\subsection{Стационарное уравнение Шрёдингера}

Нестационарное уравнение Шрёдингера имеет вид
\[
    i\hbar\frac{\partial\Psi(\vec r,t)}{\partial t}
    =\hat H\Psi(\vec r,t)
    =-\frac{\hbar^2}{2m}\Delta\Psi(\vec r,t)+U(\vec r,t)\Psi(\vec r,t),
\]
причём начальное состояние задаётся условием $\Psi(\vec r,0)=\Psi_0(\vec r)$.
Если потенциал не зависит от времени, $U(\vec r,t)=U(\vec r)$, можно применить
метод разделения переменных и искать решение в форме
\[
    \Psi(\vec r,t)=\psi(\vec r)e^{-\frac{i}{\hbar}Et}.
\]
После подстановки получаем стационарное уравнение Шрёдингера
\[
    \boxed{\hat H\psi(\vec r)=E\psi(\vec r)},
    \qquad
    \hat H=-\frac{\hbar^2}{2m}\Delta+U(\vec r).
\]
Величина $E$ имеет смысл энергии частицы. В одномерном случае
\[
    -\frac{\hbar^2}{2m}\psi''(x)+U(x)\psi(x)=E\psi(x),
    \qquad -\infty<x<+\infty,
\]
то есть возникает задача типа Штурма--Лиувилля.

Рассмотрим два решения
\[
    -\frac{\hbar^2}{2m}\Delta\psi_1+U\psi_1=E_1\psi_1,
    \qquad
    -\frac{\hbar^2}{2m}\Delta\psi_2^*+U\psi_2^*=E_2^*\psi_2^*,
    \qquad U(\vec r)\in\mathbb R.
\]
Умножая первое уравнение на $\psi_2^*$, второе на $\psi_1$ и вычитая,
получаем
\[
    -\frac{\hbar^2}{2m}\nabla\!\cdot\!
    \left(\psi_2^*\nabla\psi_1-\psi_1\nabla\psi_2^*\right)
    =(E_1-E_2^*)\psi_1\psi_2^*.
\]
Интегрирование по всему пространству при исчезновении граничного вклада
на бесконечности даёт
\[
    0=(E_1-E_2^*)\int_{\mathbb R^3}\psi_1(\vec r)\psi_2^*(\vec r)\,d^3r.
\]
Отсюда сразу следуют два результата. Если $\psi_1=\psi_2$, то $E_1=E_1^*$,
то есть $E_1\in\mathbb R$. Если же $E_1\neq E_2$, то
\[
    \int_{\mathbb R^3}\psi_1(\vec r)\psi_2^*(\vec r)\,d^3r
    =\langle\psi_2\mid\psi_1\rangle=0.
\]
Таким образом, энергетические уровни вещественны, а состояния с различными
энергиями ортогональны.

Для связанного состояния $\psi(\vec r)\in L_2(\mathbb R^3)$ и
\[
    \hat H\psi_n=E_n\psi_n.
\]
Числа $E_n$ являются уровнями энергии системы и образуют дискретный спектр
оператора $\hat H$. Если одному значению $E_n$ соответствует несколько
линейно независимых собственных функций, уровень называется вырожденным.
Для удерживающего потенциала, который локально ограничен и возрастает
на бесконечности,
\[
    U(\vec r)\longrightarrow+\infty,
    \qquad |\vec r|\longrightarrow\infty,
\]
возникает дискретный набор связанных состояний.

Пусть $\{\psi_n\}$ образует счётный ортонормированный базис:
\[
    \langle\psi_n\mid\psi_m\rangle
    =\int_{\mathbb R^3}\psi_n^*(\vec r)\psi_m(\vec r)\,d^3r
    =\delta_{nm}.
\]
Тогда произвольное состояние из $L_2(\mathbb R^3)$ раскладывается в ряд
\[
    \psi(\vec r)=\sum_n c_n\psi_n(\vec r),
    \qquad
    c_n=\langle\psi_n\mid\psi\rangle
    =\int_{\mathbb R^3}\psi_n^*(\vec r)\psi(\vec r)\,d^3r,
\]
причём сходимость понимается по норме $L_2$. Подставляя коэффициенты обратно,
получаем
\begin{align*}
    \psi(\vec r)
    &=\sum_n\psi_n(\vec r)
      \int_{\mathbb R^3}\psi_n^*(\vec r\,')\psi(\vec r\,')\,d^3r' \\
    &=\int_{\mathbb R^3}\psi(\vec r\,')
      \left[\sum_n\psi_n^*(\vec r\,')\psi_n(\vec r)\right]d^3r'.
\end{align*}
Следовательно,
\[
    \boxed{\sum_n\psi_n^*(\vec r\,')\psi_n(\vec r)
    =\delta(\vec r-\vec r\,')},
\]
что является соотношением полноты.

Если $\hat H\psi_n=E_n\psi_n$, то $E_n$ являются возможными результатами
измерения энергии. Для состояния $\psi=\sum_n c_n\psi_n$ вероятность
получить $E_n$ равна
\[
    p_n=\frac{|c_n|^2}{\|\psi\|^2}.
\]
При нормировке $\langle\psi\mid\psi\rangle=1$ имеем
\[
    p_n=|c_n|^2,
    \qquad
    \sum_n|c_n|^2=1.
\]

Если в начальный момент
\[
    \Psi(\vec r,0)=\Psi_0(\vec r)=\sum_n c_n\psi_n(\vec r),
\]
то при независимом от времени гамильтониане
\[
    \Psi(\vec r,t)
    =\sum_n c_n\psi_n(\vec r)e^{-\frac{i}{\hbar}E_nt}
    =\sum_n c_n(t)\psi_n(\vec r),
    \qquad
    c_n(t)=c_n(0)e^{-\frac{i}{\hbar}E_nt}.
\]
Поэтому
\[
    |c_n(t)|^2=|c_n(0)|^2,
\]
то есть распределение вероятностей по энергиям не зависит от времени.
В частности, если $\Psi(\vec r,0)=\psi_n(\vec r)$, то $p_n=1$.

% ---------------------------------------------------------
\subsection{Свободное движение}
% ---------------------------------------------------------

Для свободной частицы $U(\vec r,t)=0$, поэтому
\[
    -\frac{\hbar^2}{2m}\Delta\psi=E\psi.
\]
В одномерном случае
\[
    -\frac{\hbar^2}{2m}\psi''(x)=E\psi(x).
\]
Решение можно искать в виде плоской волны
\[
    \psi_p(x)=C_p e^{\frac{i}{\hbar}px},
    \qquad
    E=\frac{p^2}{2m}.
\]
Следовательно,
\[
    E\geq0,
    \qquad
    p=\pm\sqrt{2mE}.
\]
Для $E>0$ в одномерном случае одной энергии соответствуют два направления
импульса.

Плотность потока вероятности для плоской волны равна
\[
    j=\frac{i\hbar}{2m}
    \left(\psi_p\frac{d\psi_p^*}{dx}-\psi_p^*\frac{d\psi_p}{dx}\right)
    =|C_p|^2\frac{p}{m}.
\]
В трёхмерном случае
\[
    \vec j=|C_{\vec p}|^2\frac{\vec p}{m}.
\]
Так как
\[
    -i\hbar\frac{d}{dx}\psi_p=p\psi_p,
\]
оператор импульса имеет вид
\[
    \boxed{\hat p=-i\hbar\frac{d}{dx}},
    \qquad
    \boxed{\hat{\vec p}=-i\hbar\nabla}.
\]

Произвольную волновую функцию можно представить как суперпозицию плоских
волн:
\[
    \psi(x)=\int_{-\infty}^{+\infty}\varphi(p)C_p
    e^{\frac{i}{\hbar}px}\,dp,
\]
а коэффициент разложения определяется обратным преобразованием
\[
    \varphi(p)=\int_{-\infty}^{+\infty}\psi(x)C_p^*
    e^{-\frac{i}{\hbar}px}\,dx
    =\langle\psi_p\mid\psi\rangle.
\]
Первоначально эти формулы можно рассматривать на
$L_1(\mathbb R)\cap L_2(\mathbb R)$, после чего продолжить на $L_2(\mathbb R)$.

Из полноты плоских волн следует
\[
    \int_{-\infty}^{+\infty}\psi_p^*(x')\psi_p(x)\,dp
    =\delta(x-x').
\]
Для $\psi_p(x)=C_p e^{ipx/\hbar}$ имеем
\[
    \int_{-\infty}^{+\infty}\psi_p^*(x')\psi_p(x)\,dp
    =|C_p|^2\int_{-\infty}^{+\infty}e^{\frac{i}{\hbar}p(x-x')}\,dp
    =|C_p|^2(2\pi\hbar)\delta(x-x').
\]
Отсюда выбираем
\[
    \boxed{C_p=\frac{1}{\sqrt{2\pi\hbar}}},
    \qquad
    \psi_p(x)=\frac{1}{\sqrt{2\pi\hbar}}e^{\frac{i}{\hbar}px}.
\]
При такой нормировке
\[
    \int_{-\infty}^{+\infty}\psi_p^*(x)\psi_{p'}(x)\,dx
    =\delta(p-p'),
    \qquad
    \langle\psi_p\mid\psi_{p'}\rangle=\delta(p-p').
\]
Это непрерывный аналог ортогональности
$\langle\psi_n\mid\psi_m\rangle=\delta_{nm}$, а $|\varphi(p)|^2$ задаёт
плотность распределения импульса.

В трёхмерном случае
\[
    \psi_{\vec p}(\vec r)
    =\frac{1}{(2\pi\hbar)^{3/2}}
    e^{\frac{i}{\hbar}\vec p\cdot\vec r},
    \qquad
    p^2=2mE.
\]
При фиксированной энергии направление $\vec p$ может быть произвольным,
поэтому возникает бесконечная кратность вырождения. Импульс, скорость и волновой вектор связаны соотношениями
$\vec p=m\vec v$ и $\vec p=\hbar\vec k$. Для свободной частицы
связанных состояний нет.

Преобразование Фурье сохраняет норму:
\[
    \int_{-\infty}^{+\infty}|\psi(x)|^2\,dx
    =\int_{-\infty}^{+\infty}|\varphi(p)|^2\,dp=1
\]
для нормированного состояния. Аналогично для дискретного базиса
\[
    \langle\psi\mid\psi\rangle
    =\sum_n|c_n|^2
    =\sum_n|\langle\psi_n\mid\psi\rangle|^2,
\]
а для непрерывного импульсного базиса
\[
    \langle\psi\mid\psi\rangle
    =\int_{-\infty}^{+\infty}|\langle\psi_p\mid\psi\rangle|^2\,dp.
\]

Координатное и импульсное распределения связаны преобразованием Фурье:
\[
    f(x)=\int_{-\infty}^{+\infty}\varphi(p)
    \frac{1}{\sqrt{2\pi\hbar}}e^{\frac{i}{\hbar}px}\,dp.
\]
Если характерные ширины распределений обозначить через $\delta x$ и
$\delta p$, то их произведение имеет порядок
\[
    \delta x\,\delta p\sim2\pi\hbar.
\]
Поэтому координатное и импульсное распределения не могут быть одновременно
сколь угодно узкими; это выражает принцип неопределённости Гейзенберга.

Начальный волновой пакет запишем как
\[
    \Psi(x,0)=\int_{-\infty}^{+\infty}\varphi(p)\psi_p(x)\,dp.
\]
Его эволюция свободной частицы имеет вид
\[
    \Psi(x,t)=\int_{-\infty}^{+\infty}\varphi(p)\psi_p(x)
    e^{-\frac{i}{\hbar}E_pt}\,dp,
    \qquad
    E_p=\frac{p^2}{2m}.
\]
Сдвиг импульсного распределения на $p_0$ соответствует умножению функции
на плоскую волну:
\[
    f(x-x_0)e^{\frac{i}{\hbar}p_0(x-x_0)}
    =\int_{-\infty}^{+\infty}\varphi(p-p_0)
    \frac{1}{\sqrt{2\pi\hbar}}
    e^{\frac{i}{\hbar}p(x-x_0)}\,dp.
\]
Эквивалентно, импульсная амплитуда такого пакета имеет фазовый множитель
\[
    \widetilde\varphi(p)=\varphi(p-p_0)e^{-\frac{i}{\hbar}px_0}.
\]

Для пакета, сосредоточенного около $p_0$, имеем
\[
    \Psi(x,t)=\int_{-\infty}^{+\infty}\varphi(p-p_0)
    \frac{1}{\sqrt{2\pi\hbar}}
    \exp\!\left\{\frac{i}{\hbar}[p(x-x_0)-E_pt]\right\}\,dp.
\]
При $p\approx p_0$ фазу можно разложить по отклонению $p-p_0$:
\[
    p(x-x_0)-E_pt
    \simeq p_0(x-x_0)-E_{p_0}t
    +(p-p_0)\left(x-x_0-
    \left.\frac{dE}{dp}\right|_{p_0}t\right).
\]
Поэтому
\[
    \Psi(x,t)\simeq
    e^{\frac{i}{\hbar}[p_0(x-x_0)-E_{p_0}t]}
    f(x-x_0-v_0t),
    \qquad
    v_0=\left.\frac{dE}{dp}\right|_{p_0}=\frac{p_0}{m}.
\]
Для свободной частицы полезно также точное равенство
\[
    E(p)=\frac{p_0^2}{2m}+\frac{p_0}{m}(p-p_0)
    +\frac{(p-p_0)^2}{2m}.
\]
Линейная часть определяет перенос центра пакета, а квадратичный член
становится существенным, если распределение по импульсу уже нельзя считать
очень узким.

Нормировку по импульсу можно заменить нормировкой по другой непрерывной
переменной $f=f(p)$. Если
\[
    \int\psi_p^*(x)\psi_{p'}(x)\,dx=\delta(p-p'),
\]
то из
\[
    \delta(f-f')\,df=\delta(p-p')\,dp,
    \qquad
    \delta(f-f')=\frac{\delta(p-p')}{\left|df/dp\right|}
\]
получаем
\[
    \boxed{\psi_f(x)=\frac{1}{\sqrt{|df/dp|}}\,\psi_p(x)}.
\]
Для энергии $E=p^2/(2m)$ необходимо учитывать две ветви
$p=\alpha\sqrt{2mE}$, $\alpha=\pm1$. Тогда
\[
    \int_{-\infty}^{+\infty}dp
    =\sum_{\alpha=\pm1}\int_0^{+\infty}
      \left|\frac{dp}{dE}\right|dE,
\]
и энергетически нормированные функции можно записать как
\[
    \boxed{
    \psi_{E,\alpha}(x)=
    \frac{1}{\sqrt{2\pi\hbar v}}
    \exp\left(\frac{i}{\hbar}\alpha\sqrt{2mE}\,x\right)},
    \qquad v=\frac{|p|}{m},
\]
с нормировкой
\[
    \int_{-\infty}^{+\infty}\psi_{E,\alpha}^*(x)
    \psi_{E',\alpha'}(x)\,dx
    =\delta_{\alpha\alpha'}\delta(E-E').
\]

Чтобы получить непрерывный спектр как предел дискретного, рассмотрим
одномерный ящик длины $L_x$ с периодическими граничными условиями
\[
    \psi(0)=\psi(L_x).
\]
Для плоской волны это означает
\[
    e^{\frac{i}{\hbar}pL_x}=1,
    \qquad
    \frac{pL_x}{\hbar}=2\pi n,
    \qquad n\in\mathbb Z,
\]
поэтому
\[
    p_n=\frac{2\pi\hbar}{L_x}n,
    \qquad
    \Delta p=p_{n+1}-p_n=\frac{2\pi\hbar}{L_x}.
\]
Скалярное произведение берётся на $[0,L_x]$,
\[
    \langle\varphi\mid\psi\rangle
    =\int_0^{L_x}\varphi^*(x)\psi(x)\,dx,
\]
а нормированные функции равны
\[
    \psi_n(x)=\frac{1}{\sqrt{L_x}}e^{\frac{i}{\hbar}p_nx}.
\]
При $L_x\to\infty$ шаг $\Delta p$ стремится к нулю, и сумма переходит в
интеграл:
\[
    \boxed{\sum_p\longrightarrow\frac{L_x}{2\pi\hbar}\int dp}.
\]
В трёхмерном ящике объёма $V=L_xL_yL_z$ компоненты импульса также
дискретны:
\[
    p_x=\frac{2\pi\hbar}{L_x}n_x,
    \qquad
    p_y=\frac{2\pi\hbar}{L_y}n_y,
    \qquad
    p_z=\frac{2\pi\hbar}{L_z}n_z.
\]
Нормированная плоская волна и переход от суммы к интегралу имеют вид
\[
    \psi_{\vec p}(\vec r)=\frac{1}{\sqrt V}
    e^{\frac{i}{\hbar}\vec p\cdot\vec r},
    \qquad
    \boxed{\sum_{\vec p}\longrightarrow
    \frac{V}{(2\pi\hbar)^3}\int d^3p}.
\]

% ---------------------------------------------------------
\subsection{Наблюдаемые в квантовой механике. Средние значения}
% ---------------------------------------------------------

Физической величине $f\in\mathbb R$ сопоставляется линейный
самосопряжённый оператор $\hat f$, действующий в пространстве состояний.
Для такого оператора
\[
    \langle\varphi\mid\hat f\psi\rangle
    =\langle\psi\mid\hat f\varphi\rangle^*.
\]
В координатном представлении это можно записать как
\[
    \int d^3r\,\varphi^*(\vec r)\hat f\psi(\vec r)
    =\left(\int d^3r\,\psi^*(\vec r)\hat f\varphi(\vec r)\right)^*.
\]
При работе с дифференциальными операторами важно также учитывать область,
на которой оператор определён: формального совпадения выражений для
оператора и сопряжённого оператора недостаточно без соответствующих
граничных условий.

Если
\[
    \hat f\psi_n=f_n\psi_n,
\]
то из самосопряжённости следует $f_n=f_n^*$, то есть собственные значения
наблюдаемой вещественны. Для дискретного спектра собственные функции можно
нормировать так, что
\[
    \langle\psi_n\mid\psi_m\rangle=\delta_{nm}.
\]
Для непрерывного спектра
\[
    \langle\psi_f\mid\psi_{f'}\rangle=\delta(f-f'),
\]
а дискретные и непрерывные собственные состояния взаимно ортогональны:
\[
    \langle\psi_n\mid\psi_f\rangle=0.
\]

Если спектр содержит обе части, состояние раскладывается в виде
\[
    \psi(\vec r)=\sum_n c_n\psi_n(\vec r)
    +\int_{\operatorname{cont}(f)}df\,\varphi(f)\psi_f(\vec r),
\]
и его норма равна
\[
    \langle\psi\mid\psi\rangle
    =\sum_n|c_n|^2
    +\int_{\operatorname{cont}(f)}|\varphi(f)|^2\,df.
\]
Для нормированного состояния эта величина равна единице.

Возможными результатами измерения $f$ являются значения из спектра
оператора $\hat f$. Для дискретного значения
\[
    p(f_n)=|c_n|^2=|\langle\psi_n\mid\psi\rangle|^2,
\]
а в непрерывной части
\[
    dp(f)=|\varphi(f)|^2df
    =|\langle\psi_f\mid\psi\rangle|^2df.
\]
Если система находится в стационарном состоянии
\[
    \Psi(\vec r,t)=\psi(\vec r)e^{-\frac{i}{\hbar}Et},
\]
то общий фазовый множитель не меняет вероятностное распределение.

Среднее значение наблюдаемой равно
\[
    \overline f
    =\sum_n f_n|c_n|^2
    +\int_{\operatorname{cont}(f)}f|\varphi(f)|^2\,df
    =\boxed{\langle\psi\mid\hat f\psi\rangle}.
\]
Разброс результатов измерения характеризуется величиной
\[
    \operatorname{Var}f
    =\overline{(f-\overline f)^2}.
\]

Оператор координаты действует умножением:
\[
    \hat x\psi=x\psi,
\]
а оператор импульса в координатном представлении имеет вид
\[
    \hat{\vec p}\psi=-i\hbar\nabla\psi.
\]

\medskip
\noindent\textbf{Глава 2. Волновые задачи квантовой механики.}

% ---------------------------------------------------------
\subsection{Прямоугольные потенциальные ямы (1D)}
% ---------------------------------------------------------

\subsubsection{Яма с бесконечными стенками}

Рассмотрим потенциал
\[
    U(x)=
    \begin{cases}
        0, & 0<x<a,\\
        +\infty, & \text{иначе}.
    \end{cases}
\]
Внутри ямы стационарное уравнение имеет вид
\[
    -\frac{\hbar^2}{2m}\psi''=E\psi,
\]
а бесконечные стенки задают граничные условия
\[
    \psi(0)=\psi(a)=0.
\]
Общее решение внутри можно записать как
\[
    \psi=C_1e^{\frac{i}{\hbar}px}+C_2e^{-\frac{i}{\hbar}px},
    \qquad
    \frac{p^2}{2m}=E.
\]
После выполнения граничных условий остаётся синусоидальное решение
\[
    \psi(x)=C\sin\frac{px}{\hbar},
\]
причём
\[
    \frac{p_na}{\hbar}=\pi n,
    \qquad
    p_n=\frac{\pi\hbar}{a}n,
    \qquad n=1,2,\ldots
\]
Следовательно,
\[
    \boxed{E_n=\frac{\pi^2\hbar^2}{2ma^2}n^2},
\]
а нормированные собственные функции равны
\[
    \boxed{\psi_n(x)=\sqrt{\frac{2}{a}}\sin\frac{\pi nx}{a}},
    \qquad n=1,2,\ldots
\]

% =========================================================
\section{Лекции 5--6}
% =========================================================

\subsection{Прямоугольные потенциальные ямы (1D)}

В начале продолжим результат для бесконечно глубокой ямы. При
$k=p/\hbar$ граничное условие даёт
\[
    ka=\pi n,
    \qquad n=1,2,\ldots,
    \qquad
    \hbar k=\sqrt{2mE},
\]
поэтому
\[
    E_n=\frac{\hbar^2\pi^2n^2}{2ma^2},
    \qquad
    \psi_n(x)=\sqrt{\frac{2}{a}}\sin\frac{\pi n}{a}x.
\]

\subsubsection{Прямоугольная потенциальная яма конечной глубины}

Пусть
\[
    U(x)=
    \begin{cases}
        -U_0, & |x|<\dfrac a2,\\
        0, & |x|>\dfrac a2,
    \end{cases}
    \qquad U_0>0,
\]
и рассматривается связанное состояние с $E<0$. Удобно разделить ось на
три области: I --- $x<-a/2$, II --- $|x|<a/2$, III --- $x>a/2$.

В областях I и III
\[
    -\frac{\hbar^2}{2m}\psi''=E\psi,
    \qquad
    \hbar\chi=\sqrt{2m|E|},
\]
поэтому решения имеют вид $e^{\pm\chi x}$. Нормируемость оставляет только
затухающие экспоненты на соответствующих бесконечностях.

В области II
\[
    -\frac{\hbar^2}{2m}\psi''=(E+U_0)\psi,
    \qquad
    \hbar k=\sqrt{2m(E+U_0)}.
\]
Введём также
\[
    \hbar k_0=\sqrt{2mU_0}.
\]
Тогда параметры связаны соотношением
\[
    \boxed{k^2+\chi^2=k_0^2}.
\]
На границах ямы сшиваются сама волновая функция и её производная:
\[
    \psi_{\mathrm I}\!\left(-\frac a2\right)
    =\psi_{\mathrm{II}}\!\left(-\frac a2\right),
    \qquad
    \psi'_{\mathrm I}\!\left(-\frac a2\right)
    =\psi'_{\mathrm{II}}\!\left(-\frac a2\right),
\]
и аналогично при $x=a/2$.

Число связанных уровней оценивается как
\[
    \boxed{N=\left[\frac{k_0a}{\pi}\right]+1}.
\]
В одномерной симметричной яме существует хотя бы одно связанное состояние.
Для глубокой ямы $k_0a\gg1$ уровней много, а для мелкой ямы
\[
    k_0a\ll1
\]
остаётся одно связанное состояние: $N=1$.

В мелкой яме одновременно $ka\ll1$ и $\chi a\ll1$. Основное состояние
чётно, поэтому удобно взять
\[
    \psi_{\mathrm I}=C_1e^{\chi x},
    \qquad
    \psi_{\mathrm{II}}=A\cos kx,
    \qquad
    \psi_{\mathrm{III}}=C_1e^{-\chi x}.
\]
Условия сшивки при $x=a/2$ имеют вид
\[
    A\cos\frac{ka}{2}=C_1e^{-\chi a/2},
    \qquad
    -Ak\sin\frac{ka}{2}=-C_1\chi e^{-\chi a/2}.
\]
При малых $ka$ и $\chi a$ отсюда $A\simeq C_1$ и
\[
    k\tan\frac{ka}{2}=\chi
    \quad\Rightarrow\quad
    \chi\simeq k\frac{ka}{2}.
\]
Так как $\chi\ll k$, имеем $k\simeq k_0$ и
\[
    \chi\simeq k_0\frac{k_0a}{2}.
\]
Энергия лежит очень близко к границе $E=0$:
\[
    |E|=\frac{\hbar^2\chi^2}{2m}
    \simeq U_0\left(\frac{k_0a}{2}\right)^2\ll U_0,
    \qquad E<0.
\]

Вероятность обнаружить частицу снаружи ямы оценивается как
\[
    P_{\mathrm{out}}
    =2\int_{a/2}^{\infty}C_1^2e^{-2\chi x}\,dx
    =\frac{C_1^2}{\chi}e^{-\chi a},
\]
а внутри
\[
    P_{\mathrm{in}}
    =\int_{-a/2}^{a/2}A^2\cos^2(kx)\,dx
    \simeq aA^2.
\]
При $A\simeq C_1$ и $\chi a\ll1$
\[
    \frac{P_{\mathrm{in}}}{P_{\mathrm{out}}}
    \simeq\chi a
    =\frac{(k_0a)^2}{2}\ll1.
\]
То есть у очень мелкой ямы связанное состояние может быть заметно
растянуто за пределы самой ямы.

% ---------------------------------------------------------
\subsection{Прямоугольный потенциальный барьер}
% ---------------------------------------------------------

Рассмотрим барьер
\[
    U(x)=
    \begin{cases}
        0, & x<0,\\
        U_0, & 0<x<a,\\
        0, & x>a.
    \end{cases}
\]
Для начала пусть $0<E<U_0$. Введём
\[
    \hbar k=\sqrt{2mE},
    \qquad
    \hbar\chi=\sqrt{2m(U_0-E)},
    \qquad
    \hbar k_0=\sqrt{2mU_0},
\]
так что
\[
    k^2+\chi^2=k_0^2.
\]
Для волны, падающей слева,
\[
    \psi_k^{(L)}(x)=
    \begin{cases}
        e^{ikx}+r e^{-ikx}, & x<0,\\
        Ae^{\chi x}+Be^{-\chi x}, & 0<x<a,\\
        de^{ikx}, & x>a.
    \end{cases}
\]
Здесь $r$ и $d$ --- амплитуды отражённой и прошедшей волн.

Плотность потока вероятности равна
\[
    j=\frac{i\hbar}{2m}
    \left(\psi\frac{d\psi^*}{dx}-\psi^*\frac{d\psi}{dx}\right).
\]
Для прошедшей волны
\[
    j_{\mathrm{III}}=|d|^2\frac{\hbar k}{m},
\]
а слева
\[
    j_{\mathrm I}=\frac{\hbar k}{m}
    -|r|^2\frac{\hbar k}{m}.
\]
Из $j_{\mathrm I}=j_{\mathrm{III}}$ следует
\[
    \boxed{1=|r|^2+|d|^2=R+D},
\]
где
\[
    R=|r|^2,
    \qquad
    D=|d|^2
\]
--- коэффициенты отражения и прохождения.

Для волны, падающей справа, можно аналогично записать
\[
    \psi_k^{(R)}(x)=
    \begin{cases}
        d'e^{-ikx}, & x<0,\\
        A'e^{\chi x}+B'e^{-\chi x}, & 0<x<a,\\
        e^{-ikx}+r'e^{ikx}, & x>a.
    \end{cases}
\]

Для обсуждения движения локализованного состояния рассмотрим волновой
пакет. Если
\[
    f(x)=\frac{1}{\sqrt{2\pi}}\int\varphi(k)e^{ikx}\,dk,
\]
то сдвиг и умножение на несущую волну дают соответствие
\[
    f(x-x_0)e^{ik_0(x-x_0)}
    \longleftrightarrow
    \varphi(k-k_0)e^{-ikx_0}.
\]
Для пакета, падающего слева,
\[
    \Psi(x,t)=\int_0^\infty dk\,\varphi(k-k_0)e^{-ikx_0}
    \psi_k^{(L)}(x)e^{-\frac{i}{\hbar}E_kt}.
\]
В области $x<0$ присутствуют падающая и отражённая части. Запишем
\[
    r(k)=|r(k)|e^{i\alpha_r(k)},
    \qquad
    d(k)=|d(k)|e^{i\alpha_d(k)}.
\]
Для узкого пакета фазы можно разложить около $k_0$. Например, фаза
отражённой компоненты содержит
\[
    -k(x+x_0)-\frac{E_k}{\hbar}t+\alpha_r(k),
\]
поэтому в линейном приближении
\begin{align*}
    -k(x+x_0)-\frac{E_k}{\hbar}t+\alpha_r(k)
    &\simeq -k_0(x+x_0)-\frac{E_{k_0}}{\hbar}t+\alpha_r(k_0) \\
    &\quad +(k-k_0)\left[-x-x_0-\frac{1}{\hbar}
    \left.\frac{dE}{dk}\right|_{k_0}t
    +\left.\frac{d\alpha_r}{dk}\right|_{k_0}\right].
\end{align*}
Так как $v_0=(1/\hbar)(dE/dk)_{k_0}=\hbar k_0/m$, отражённая часть
пакета имеет форму
\[
    \Psi_{\mathrm{ref}}(x,t)
    \sim |r(k_0)|e^{i\Phi_r}
    f\bigl(-x-x_0-v_0(t-\tau_r)\bigr),
\]
а её центр движется как
\[
    x=-x_0-v_0(t-\tau_r),
    \qquad
    \boxed{\tau_r=\frac{1}{v_0}
    \left.\frac{d\alpha_r}{dk}\right|_{k_0}}.
\]
Для прошедшего пакета аналогично
\[
    \Psi_{\mathrm{tr}}(x,t)
    \sim |d(k_0)|e^{i\Phi_d}
    f\bigl(x-x_0-v_0(t-\tau_d)\bigr),
\]
\[
    \boxed{\tau_d=\frac{1}{v_0}
    \left.\frac{d\alpha_d}{dk}\right|_{k_0}}.
\]
Поэтому момент прохождения точки $x=a$ можно записать как
\[
    t=\frac{a-x_0}{v_0}+\tau_d.
\]

Для $E<U_0$ коэффициент прохождения прямоугольного барьера равен
\[
    \boxed{
    D=\frac{1}{1+\left(\dfrac{k_0^2}{2k\chi}\right)^2
    \sinh^2(\chi a)}}.
\]
Если $\chi a\ll1$, то $\sinh(\chi a)\simeq\chi a$, и
\[
    D\simeq
    \frac{1}{1+\left(\dfrac{k_0}{k}\right)^2
    \left(\dfrac{k_0a}{2}\right)^2}
    =\frac{1}{1+\dfrac{U_0}{E}
    \left(\dfrac{k_0a}{2}\right)^2}.
\]
При $E=U_0$
\[
    \boxed{D(E=U_0)=\frac{1}{1+\left(\dfrac{k_0a}{2}\right)^2}}.
\]

Для тонкого барьера $k_0a\ll1$, следовательно, также $ka\ll1$ и
$\chi a\ll1$. При
\[
    E\ll U_0\left(\frac{k_0a}{2}\right)^2
\]
получаем низкоэнергетическую оценку
\[
    D\simeq\frac{E}{U_0}
    \left(\frac{k_0a}{2}\right)^{-2},
    \qquad k\to0.
\]
Если $k\to0$, то $\chi\simeq k_0$, и из общей формулы следует
\[
    D\simeq4\frac{E}{U_0}\,\sinh^{-2}(k_0a).
\]
Для толстого барьера $k_0a\gg1$ главным является экспоненциальное
подавление
\[
    D\sim\frac{E}{U_0}e^{-2k_0a}.
\]

При $E>U_0$ внутри барьера решение снова становится осциллирующим.
Положим
\[
    \chi\to ik_2,
    \qquad
    \hbar k_2=\sqrt{2m(E-U_0)}.
\]
Тогда
\[
    \boxed{
    D=\frac{1}{1+\left(\dfrac{k_0^2}{2kk_2}\right)^2
    \sin^2(k_2a)}}.
\]
Если
\[
    k_2a=\pi n,
    \qquad n=1,2,\ldots,
\]
то
\[
    \boxed{D=1,\qquad R=0}.
\]
Это резонансные значения энергии: барьер становится полностью прозрачным.
Между резонансами минимальное значение равно
\[
    \boxed{D_{\min}=\frac{1}{1+
    \left(\dfrac{k_0^2}{2kk_2}\right)^2}}.
\]

% ---------------------------------------------------------
\subsection{Линейный гармонический осциллятор}
% ---------------------------------------------------------

Для линейного гармонического осциллятора
\[
    U(x)=\frac{m\omega^2x^2}{2},
\]
поэтому стационарное уравнение Шрёдингера имеет вид
\[
    -\frac{\hbar^2}{2m}\psi''
    +\frac{m\omega^2x^2}{2}\psi=E\psi.
\]
Решение удобно строить в три шага: перейти к безразмерным переменным,
исследовать асимптотику и затем искать оставшуюся часть решения в виде ряда.

Введём
\[
    x=a\xi,
    \qquad
    a=\sqrt{\frac{\hbar}{m\omega}},
    \qquad
    E=\frac{\hbar\omega}{2}\lambda.
\]
Тогда уравнение принимает вид
\[
    \boxed{\psi''+(\lambda-\xi^2)\psi=0},
\]
где штрих означает производную по $\xi$.

Для исследования асимптотики введём
\[
    Y=\frac{\psi'}{\psi}.
\]
Тогда
\[
    Y'=\frac{\psi''}{\psi}
    -\left(\frac{\psi'}{\psi}\right)^2
    =\frac{\psi''}{\psi}-Y^2,
\]
и исходное уравнение переписывается как
\[
    Y'+Y^2+\lambda-\xi^2=0.
\]
При $|\xi|\to\infty$ доминирует $\xi^2$, поэтому
\[
    Y\simeq\pm\xi,
    \qquad
    \frac{\psi'}{\psi}\simeq\pm\xi.
\]
Отсюда
\[
    \ln\psi\simeq\pm\frac{\xi^2}{2}+C,
    \qquad
    \psi\sim e^{\pm\xi^2/2}.
\]
Растущее решение не нормируется, поэтому оставляем затухающую асимптотику
и записываем
\[
    \psi(\xi)=e^{-\xi^2/2}H(\xi).
\]
Подстановка даёт
\[
    \boxed{H''-2\xi H'+(\lambda-1)H=0}.
\]

Ищем
\[
    H(\xi)=\sum_{n=0}^{\infty}a_n\xi^n,
    \qquad
    H'=\sum_{n=1}^{\infty}na_n\xi^{n-1},
    \qquad
    H''=\sum_{n=2}^{\infty}n(n-1)a_n\xi^{n-2}.
\]
Сравнение коэффициентов при одинаковых степенях $\xi$ приводит к
\[
    \boxed{(n+1)(n+2)a_{n+2}=(2n+1-\lambda)a_n}.
\]
Чётные коэффициенты связаны только с чётными, а нечётные --- только с
нечётными. При $n\to\infty$
\[
    a_{n+2}\simeq\frac{2}{n}a_n.
\]
Например, для чётной последовательности $n=2k$ получаем
\[
    a_{2(k+1)}\simeq\frac{1}{k}a_{2k},
    \qquad
    a_0\sum_{k=0}^{\infty}\frac{\xi^{2k}}{k!}
    =a_0e^{\xi^2}.
\]
Следовательно, бесконечный ряд приводил бы к росту $H$ порядка
$e^{\xi^2}$, то есть к ненормируемой $\psi$. Поэтому ряд должен
оборваться.

Обрезание происходит при
\[
    \boxed{\lambda=2n+1},
    \qquad n=0,1,2,\ldots
\]
Для чётных $n$ остаётся чётное решение ($a_1=0$), для нечётных --- нечётное
($a_0=0$). Спектр энергии равен
\[
    \boxed{E_n=\hbar\omega\left(n+\frac12\right)},
    \qquad
    E_{n+1}-E_n=\hbar\omega.
\]

При разрешённых $\lambda$ функция $H$ становится полиномом Эрмита $H_n$,
который удовлетворяет уравнению
\[
    \boxed{H_n''-2\xi H_n'+2nH_n=0}.
\]
Поэтому
\[
    \psi_n=C_ne^{-\xi^2/2}H_n(\xi).
\]
Для полиномов Эрмита удобно ввести производящую функцию
\[
    \boxed{\mathcal U(s,\xi)=e^{-s^2+2s\xi}
    =\sum_{n=0}^{\infty}\frac{s^n}{n!}H_n(\xi)}.
\]
Из
\[
    \frac{\partial\mathcal U}{\partial\xi}=2s\mathcal U
\]
следует
\[
    H_{n+1}'(\xi)=2(n+1)H_n(\xi),
\]
а из
\[
    \frac{\partial\mathcal U}{\partial s}=2(\xi-s)\mathcal U
\]
получаем рекуррентное соотношение
\[
    H_{n+1}(\xi)=2\xi H_n(\xi)-2nH_{n-1}(\xi).
\]
Эти равенства вновь приводят к уравнению Эрмита
\[
    H_n''-2\xi H_n'+2nH_n=0.
\]
Полиномы можно также записать по формуле Родрига:
\[
    \boxed{H_n(\xi)=(-1)^ne^{\xi^2}
    \frac{d^n}{d\xi^n}e^{-\xi^2}}.
\]
Первые из них:
\[
    H_0(\xi)=1,
    \qquad
    H_1(\xi)=2\xi,
    \qquad
    H_2(\xi)=4\xi^2-2.
\]

Нормировка определяется ортогональностью
\[
    \int_{-\infty}^{+\infty}e^{-\xi^2}H_n(\xi)H_m(\xi)\,d\xi
    =\sqrt\pi\,2^n n!\,\delta_{nm}.
\]
Это соотношение можно получить из производящей функции:
\begin{align*}
    \int_{-\infty}^{+\infty}e^{-\xi^2}
    \mathcal U(s,\xi)\mathcal U(t,\xi)\,d\xi
    &=\int_{-\infty}^{+\infty}
      e^{-\xi^2-s^2-t^2+2(s+t)\xi}\,d\xi \\
    &=\sqrt\pi\,e^{2st}
    =\sqrt\pi\sum_{n=0}^{\infty}\frac{2^n s^nt^n}{n!}.
\end{align*}
В результате
\[
    1=a|C_n|^2\sqrt\pi\,2^n n!,
    \qquad a=\sqrt{\frac{\hbar}{m\omega}},
\]
и нормированная собственная функция принимает вид
\[
    \boxed{
    \psi_n(x)=
    \left(\frac{m\omega}{\pi\hbar}\right)^{1/4}
    \frac{1}{\sqrt{2^n n!}}
    e^{-\frac{m\omega x^2}{2\hbar}}
    H_n\!\left(\sqrt{\frac{m\omega}{\hbar}}x\right)}.
\]

% ---------------------------------------------------------
\subsection{Общие свойства одномерного движения}
% ---------------------------------------------------------

Рассмотрим локально ограниченный одномерный потенциал, который на больших
$|x|$ стремится к конечным значениям $U_-$ слева и $U_+$ справа. Поведение
решений удобно разделить по положению энергии относительно этих предельных
значений и минимума $U_{\min}$.

Если
\[
    E>U_{\min},
    \qquad
    E<U_-,\;U_+,
\]
классическое движение ограничено. В квантовой задаче этому соответствует
невырожденный дискретный спектр связанных состояний; число связанных
уровней для такой конечной ямы конечно. Собственные функции упорядочиваются
по числу узлов: у $n$-го связанного состояния ровно $n$ узлов,
\[
    n=0,1,2,\ldots
\]

Если энергия превышает один из асимптотических уровней потенциала, но ещё
не превышает другой, возникает непрерывный невырожденный спектр.

Если же
\[
    E>U_-,\;U_+,
\]
движение возможно на обеих бесконечностях и возникает вырожденный
непрерывный спектр: при одной и той же энергии возможны состояния,
соответствующие движению с разных сторон.

% =========================================================
\section{Лекции 9--10}
% =========================================================

\subsection{Дираковская решётка}

Для дираковской решётки связь между волновым числом $k$ и
квазиимпульсом $q$ задаётся дисперсионным соотношением
\[
    \boxed{
        \cos(qa)
        =
        \cos(ka)+\frac{\Omega}{k}\sin(ka)
    }.
\]
Обозначим правую часть через
\[
    F(k)
    =
    \cos(ka)+\frac{\Omega}{k}\sin(ka).
\]
Тогда уравнение для разрешённых значений $k$ имеет компактный вид
\[
    \cos(qa)=F(k).
\]

Функцию $F(k)$ можно представить как одну гармоническую функцию с
зависящими от $k$ амплитудой и фазой:
\[
    F(k)
    =
    \sqrt{1+\frac{\Omega^2}{k^2}}\,
    \cos\bigl(ka+\varphi(k)\bigr).
\]
В точках $k=\pi n/a$ имеем
\[
    F\left(\frac{\pi n}{a}\right)=(-1)^n,
    \qquad
    F'\left(\frac{\pi n}{a}\right)
    =
    \frac{\Omega a}{k}(-1)^n.
\]

При малых $k$ удобно использовать разложения
\[
    \cos(ka)
    =
    1-\frac{(ka)^2}{2}+\ldots,
    \qquad
    \frac{\Omega}{k}\sin(ka)
    =
    \Omega a
    \left(
        1-\frac{(ka)^2}{6}+\ldots
    \right),
\]
так что
\[
    F(k)
    =
    1+\Omega a
    -
    \left(
        \frac12+\frac{\Omega a}{6}
    \right)(ka)^2
    +\ldots
\]

Поскольку
\[
    -1\leq \cos(qa)\leq1,
\]
вещественные значения квазиимпульса существуют только для тех $k$,
для которых
\[
    \boxed{|F(k)|\leq1}.
\]
Поэтому энергетический спектр разбивается на разрешённые зоны,
разделённые запрещёнными областями.

Энергия выражается через $k$ обычным образом:
\[
    \boxed{
        E=\frac{\hbar^2k^2}{2m}
    }.
\]
Квазиимпульс можно ограничить первым интервалом
\[
    \boxed{
        q\in
        \left[
            -\frac{\pi}{a},
            \frac{\pi}{a}
        \right]
    }.
\]
Для каждого $q$ дисперсионное уравнение имеет несколько решений,
которые удобно обозначать $k_n(q)$. Поэтому энергетические ветви
записываются как
\[
    E_n(q)
    =
    \frac{\hbar^2k_n^2(q)}{2m}.
\]
Индекс $n$ нумерует энергетические зоны.

При больших $k$
\[
    \frac{\Omega}{k}\longrightarrow0,
\]
поэтому
\[
    \boxed{
        \cos(qa)\simeq\cos(ka),
        \qquad k\to\infty
    }.
\]
Таким образом, при больших энергиях дисперсия всё больше приближается
к дисперсии свободной частицы, хотя зонная структура сохраняется.


% =========================================================
\subsection{Движение в центральном поле}
% =========================================================

Рассмотрим движение частицы массы $\mu$ в трёхмерном центральном поле.
Потенциал зависит только от расстояния до центра:
\[
    V(\vec r)=V(r),
    \qquad
    r=|\vec r|.
\]
Стационарное уравнение Шрёдингера имеет вид
\[
    \boxed{
        -\frac{\hbar^2}{2\mu}\Delta\psi
        +V(r)\psi
        =
        E\psi
    }.
\]

Для центрального поля естественно перейти к сферическим координатам:
\[
    x=r\sin\theta\cos\varphi,
    \qquad
    y=r\sin\theta\sin\varphi,
    \qquad
    z=r\cos\theta.
\]
В этих координатах оператор Лапласа равен
\[
    \Delta
    =
    \frac{1}{r^2}
    \frac{\partial}{\partial r}
    \left(
        r^2\frac{\partial}{\partial r}
    \right)
    +
    \frac{1}{r^2\sin\theta}
    \frac{\partial}{\partial\theta}
    \left(
        \sin\theta
        \frac{\partial}{\partial\theta}
    \right)
    +
    \frac{1}{r^2\sin^2\theta}
    \frac{\partial^2}{\partial\varphi^2}.
\]

Ищем решение методом разделения переменных:
\[
    \boxed{
        \psi(r,\theta,\varphi)
        =
        R(r)\Theta(\theta)\Phi(\varphi)
    }.
\]

\subsubsection{Угловая часть}

Для зависимости от азимутального угла $\varphi$ получается уравнение
\[
    \Phi''+\nu\Phi=0.
\]
Так как $\varphi$ и $\varphi+2\pi$ соответствуют одной и той же точке,
волновая функция должна удовлетворять условию периодичности
\[
    \Phi(0)=\Phi(2\pi),
    \qquad
    \varphi\in[0,2\pi].
\]
Отсюда
\[
    \nu=m^2,
    \qquad
    m\in\mathbb Z,
\]
а решение можно выбрать в виде
\[
    \boxed{
        \Phi_m(\varphi)=e^{im\varphi}
    }.
\]

Уравнение для $\Theta(\theta)$ имеет вид
\[
    \boxed{
    \frac{1}{\sin\theta}
    \frac{d}{d\theta}
    \left(
        \sin\theta\frac{d\Theta}{d\theta}
    \right)
    +
    \left(
        \lambda-\frac{m^2}{\sin^2\theta}
    \right)\Theta
    =
    0
    }.
\]
Здесь
\[
    \theta\in[0,\pi],
\]
и требуется выбрать решения, регулярные при $\theta=0$ и $\theta=\pi$.

Это условие приводит к квантованию постоянной разделения:
\[
    \boxed{
        \lambda=l(l+1)
    },
\]
где
\[
    l=|m|,\,|m|+1,\,|m|+2,\ldots
\]
или, если рассматривать все возможные $m$,
\[
    l=0,1,2,\ldots,
    \qquad
    m=-l,-l+1,\ldots,l.
\]

Решения углового уравнения выражаются через присоединённые полиномы
Лежандра:
\[
    \Theta(\theta)
    =
    P_l^{|m|}(\cos\theta),
\]
где
\[
    P_l^{|m|}(\cos\theta)
    =
    \sin^{|m|}\theta
    \left(
        \frac{d}{d(\cos\theta)}
    \right)^{|m|}
    P_l(\cos\theta).
\]
Обычные полиномы Лежандра задаются формулой
\[
    \boxed{
    P_l(x)
    =
    \frac{1}{2^l l!}
    \left(
        \frac{d}{dx}
    \right)^l
    (x^2-1)^l
    }.
\]

Полное угловое решение называется сферической функцией:
\[
    \boxed{
        Y_{lm}(\theta,\varphi)
        =
        C_{lm}
        P_l^{|m|}(\cos\theta)e^{im\varphi}
    }.
\]

Нормировочная постоянная удовлетворяет
\[
    \boxed{
    |C_{lm}|^2
    =
    \frac{2l+1}{4\pi}
    \frac{(l-|m|)!}{(l+|m|)!}
    }.
\]
При выборе фазы Кондона--Шортли можно записать
\[
    C_{lm}
    =
    (-1)^{\frac{m+|m|}{2}}
    \sqrt{
        \frac{2l+1}{4\pi}
        \frac{(l-|m|)!}{(l+|m|)!}
    }.
\]
При таком соглашении
\[
    \boxed{
        Y_{lm}^*(\theta,\varphi)
        =
        (-1)^mY_{l,-m}(\theta,\varphi)
    }.
\]

Сферические функции ортонормированы:
\[
    \boxed{
    \int d\Omega\,
    Y_{lm}(\theta,\varphi)
    Y_{l'm'}^*(\theta,\varphi)
    =
    \delta_{ll'}\delta_{mm'}
    },
\]
где
\[
    d\Omega=\sin\theta\,d\theta\,d\varphi.
\]
Поэтому произвольную функцию углов можно раскладывать как
\[
    f(\theta,\varphi)
    =
    \sum_{l=0}^{\infty}
    \sum_{m=-l}^{l}
    a_{lm}Y_{lm}(\theta,\varphi).
\]

Соотношение полноты имеет вид
\[
    \boxed{
    \sum_{l=0}^{\infty}
    \sum_{m=-l}^{l}
    Y_{lm}(\theta,\varphi)
    Y_{lm}^*(\theta',\varphi')
    =
    \frac{\delta(\theta-\theta')}{\sin\theta}
    \delta(\varphi-\varphi')
    }.
\]

\subsubsection{Момент импульса}

Оператор импульса в координатном представлении равен
\[
    \hat{\vec p}=-i\hbar\nabla.
\]
Оператор момента импульса определяется как
\[
    \boxed{
        \hat{\vec M}
        =
        [\vec r\times\hat{\vec p}]
    }.
\]
Удобно выделить множитель $\hbar$:
\[
    \hat{\vec M}=\hbar\hat{\vec L},
\]
где $\hat{\vec L}$ --- безразмерный оператор момента импульса.

Для $z$-компоненты
\[
    \hat L_z
    =
    -i
    \left(
        x\frac{\partial}{\partial y}
        -
        y\frac{\partial}{\partial x}
    \right)
    =
    -i\frac{\partial}{\partial\varphi}.
\]
Следовательно,
\[
    \boxed{
        \hat L_zY_{lm}=mY_{lm}
    }.
\]
Таким образом, $m$ является квантовым числом проекции момента импульса
на выбранную ось $z$.

Квадрат момента импульса задаётся оператором
\[
    \hat L^2
    =
    \hat L_x^2+\hat L_y^2+\hat L_z^2
    =
    -\Delta_{S^2},
\]
и сферические функции являются его собственными функциями:
\[
    \boxed{
        \hat L^2Y_{lm}
        =
        l(l+1)Y_{lm}
    }.
\]

\subsubsection{Радиальное уравнение}

После разделения угловых переменных радиальная функция удовлетворяет
уравнению
\[
    \boxed{
    -\frac{\hbar^2}{2\mu r^2}
    \frac{d}{dr}
    \left(
        r^2\frac{dR}{dr}
    \right)
    +
    \left(
        V(r)+
        \frac{\hbar^2l(l+1)}{2\mu r^2}
    \right)R
    =
    ER
    }.
\]
Эквивалентно,
\[
    -\frac{\hbar^2}{2\mu}
    \left(
        R''+\frac{2}{r}R'
    \right)
    +
    \left(
        V(r)+
        \frac{\hbar^2l(l+1)}{2\mu r^2}
    \right)R
    =
    ER.
\]

Член
\[
    \frac{\hbar^2l(l+1)}{2\mu r^2}
\]
добавляется к исходному потенциалу и входит в эффективный радиальный
потенциал:
\[
    \boxed{
        V_{\mathrm{eff}}(r)
        =
        V(r)+
        \frac{\hbar^2l(l+1)}{2\mu r^2}
    }.
\]
При этом
\[
    \hbar^2l(l+1)
\]
является собственным значением квадрата физического момента импульса
$\hat{\vec M}^{\,2}$.

Чтобы убрать первую производную, вводят
\[
    R(r)=\frac{\chi(r)}{r}.
\]
Тогда
\[
    R'
    =
    \frac{\chi'}{r}
    -
    \frac{\chi}{r^2},
    \qquad
    R''
    =
    \frac{\chi''}{r}
    -
    \frac{2\chi'}{r^2}
    +
    \frac{2\chi}{r^3},
\]
и радиальное уравнение принимает вид
\[
    \boxed{
    -\frac{\hbar^2}{2\mu}\chi''
    +
    \left(
        V(r)+
        \frac{\hbar^2l(l+1)}{2\mu r^2}
    \right)\chi
    =
    E\chi
    }.
\]
По форме это одномерное уравнение Шрёдингера на полуоси
\[
    r\in[0,\infty).
\]

Для орбитального квантового числа традиционно используются обозначения
\[
    \begin{array}{c|cccccc}
        l & 0 & 1 & 2 & 3 & 4 & 5 \\ \hline
          & s & p & d & f & g & h
    \end{array}.
\]

Если $V(r)<0$ и потенциал стремится к нулю при $r\to\infty$, возможны
связанные состояния с $E<0$ и состояния непрерывного спектра с $E>0$.
Радиальные функции соответственно можно обозначать
\[
    R_{nl}(r)
    \quad\text{для связанных состояний},
    \qquad
    R_{El}(r)
    \quad\text{для непрерывного спектра}.
\]

Рассмотрим поведение решения при $r\to0$. Если сингулярность $V(r)$
слабее, чем $1/r^2$, главным становится центробежный член. Полагая
\[
    R(r)\sim r^s,
\]
получаем
\[
    s(s-1)+2s-l(l+1)=0,
\]
то есть
\[
    s(s+1)=l(l+1).
\]
Следовательно,
\[
    s=l
    \qquad\text{или}\qquad
    s=-l-1.
\]
Регулярному решению соответствует
\[
    \boxed{
        R(r)\sim r^l,
        \qquad
        \chi(r)=rR(r)\sim r^{l+1},
        \qquad
        r\to0
    }.
\]


% =========================================================
\subsubsection{Пример: свободное движение, \texorpdfstring{$V=0$}{V=0}}
% =========================================================

При $V=0$ энергия непрерывного спектра записывается как
\[
    \boxed{
        E=\frac{\hbar^2k^2}{2\mu}
    }.
\]
Для $l=0$ уравнение для $\chi$ принимает особенно простой вид:
\[
    \chi''+k^2\chi=0.
\]
Регулярное в начале координат решение можно выбрать как
\[
    \chi(r)=\sin(kr),
\]
поэтому
\[
    R_{k0}(r)\propto\frac{\sin(kr)}{r}.
\]

При $l\neq0$ радиальное уравнение имеет вид
\[
    R''+\frac{2}{r}R'
    +
    \left(
        k^2-\frac{l(l+1)}{r^2}
    \right)R
    =
    0.
\]
После введения
\[
    x=kr
\]
регулярные решения выражаются через сферические функции Бесселя
$j_l(x)$:
\[
    \boxed{
    j_l(x)
    =
    (-1)^l x^l
    \left(
        \frac{1}{x}\frac{d}{dx}
    \right)^l
    \frac{\sin x}{x}
    }.
\]

При больших $x$
\[
    \boxed{
        j_l(x)
        \sim
        \frac{
            \sin\left(x-\dfrac{\pi l}{2}\right)
        }{x}
    }.
\]
Если потенциал $V(r)$ отличен от нуля, но стремится к нулю на
бесконечности, асимптотика приобретает дополнительный фазовый сдвиг:
\[
    \boxed{
        R_{kl}(r)
        \sim
        C_{kl}
        \frac{
            \sin\left(
                kr-\dfrac{\pi l}{2}+\delta_l
            \right)
        }{kr}
    },
    \qquad
    r\to\infty.
\]

Состояния непрерывного спектра нормируются на $\delta$-функцию:
\[
    \boxed{
    \int_0^\infty
        r^2\,dr\,
        R_{kl}(r)R_{k'l}^*(r)
    =
    \delta(k-k')
    }.
\]
Используя асимптотический вид,
\[
    \int_0^\infty
    \sin\left(kr-\frac{\pi l}{2}\right)
    \sin\left(k'r-\frac{\pi l}{2}\right)\,dr
    =
    \frac{\pi}{2}\delta(k-k'),
\]
получаем
\[
    \frac{|C_{kl}|^2}{kk'}\,
    \frac{\pi}{2}\delta(k-k')
    =
    \delta(k-k').
\]
Следовательно, можно выбрать
\[
    \boxed{
        C_{kl}
        =
        \sqrt{\frac{2}{\pi}}\,k
    }.
\]


% =========================================================
\subsection{Водородоподобный ион}
% =========================================================

Рассмотрим систему из ядра массы $m_A$ и электрона массы $m_e$.
Гамильтониан двух частиц имеет вид
\[
    \boxed{
    \hat H
    =
    -\frac{\hbar^2}{2m_A}\Delta_A
    -
    \frac{\hbar^2}{2m_e}\Delta_e
    +
    V(\vec r_A-\vec r_e)
    }.
\]

Введём относительную координату
\[
    \boxed{
        \vec r=\vec r_A-\vec r_e
    }
\]
и координату центра масс
\[
    \boxed{
        \vec R
        =
        \frac{
            m_e\vec r_e+m_A\vec r_A
        }{
            m_e+m_A
        }
    }.
\]
Якобиан этого линейного преобразования равен единице.

Обратное преобразование имеет вид
\[
    \vec r_A
    =
    \vec R+
    \frac{m_e}{m_e+m_A}\vec r,
    \qquad
    \vec r_e
    =
    \vec R-
    \frac{m_A}{m_e+m_A}\vec r.
\]
Поэтому
\[
    \nabla_A
    =
    \nabla_r+
    \frac{m_A}{m_e+m_A}\nabla_R,
    \qquad
    \nabla_e
    =
    -\nabla_r+
    \frac{m_e}{m_e+m_A}\nabla_R.
\]

Для лапласианов получаем
\[
    \Delta_A
    =
    \Delta_r
    +
    \frac{m_A^2}{(m_e+m_A)^2}\Delta_R
    +
    2\frac{m_A}{m_e+m_A}
    \nabla_r\!\cdot\nabla_R,
\]
\[
    \Delta_e
    =
    \Delta_r
    +
    \frac{m_e^2}{(m_e+m_A)^2}\Delta_R
    -
    2\frac{m_e}{m_e+m_A}
    \nabla_r\!\cdot\nabla_R.
\]
При сложении кинетических энергий смешанные производные сокращаются.

Введём полную и приведённую массы:
\[
    \boxed{
        M=m_e+m_A,
        \qquad
        \mu=
        \frac{m_em_A}{m_e+m_A}
    }.
\]
Тогда гамильтониан разделяется:
\[
    \boxed{
    \hat H
    =
    -\frac{\hbar^2}{2M}\Delta_R
    -
    \frac{\hbar^2}{2\mu}\Delta_r
    +
    V(r)
    }.
\]

Волновая функция может быть представлена как произведение внутренней
волновой функции и плоской волны центра масс:
\[
    \boxed{
    \Psi(\vec r,\vec R)
    =
    \psi(\vec r)
    \frac{1}{(2\pi)^{3/2}}
    e^{i\vec K\cdot\vec R}
    }.
\]
Полная энергия соответственно раскладывается в сумму
\[
    \boxed{
        E
        =
        E_{\mathrm{int}}
        +
        \frac{\hbar^2K^2}{2M}
    }.
\]

Для водородоподобного иона кулоновский потенциал равен
\[
    \boxed{
        V(r)=-\frac{Ze^2}{r}
    }.
\]
Здесь $Z$ --- заряд ядра в единицах элементарного заряда.

\subsubsection{Радиальное уравнение}

Для заданного $l$ радиальное уравнение принимает вид
\[
    \boxed{
    -\frac{\hbar^2}{2\mu}
    \left(
        R''+\frac{2}{r}R'
    \right)
    +
    \left(
        -\frac{Ze^2}{r}
        +
        \frac{\hbar^2l(l+1)}{2\mu r^2}
    \right)R
    =
    ER
    }.
\]

Введём характерный радиус
\[
    \boxed{
        a
        =
        \frac{\hbar^2}{\mu e^2Z}
        =
        \frac{a_B}{Z}
    },
\]
где $a_B$ соответствует боровскому радиусу с учётом используемой
приведённой массы.

Для связанного состояния $E<0$ параметризуем энергию как
\[
    \boxed{
        E
        =
        -\frac{Ze^2}{2a\lambda^2}
    }.
\]
Здесь $\lambda$ --- безразмерный параметр, который далее будет
квантоваться.

Введём безразмерную переменную
\[
    \boxed{
        r=\frac{a\lambda}{2}\rho
    }.
\]
После подстановки радиальное уравнение принимает вид
\[
    \boxed{
    R''+
    \frac{2}{\rho}R'
    +
    \left(
        -\frac14
        -\frac{l(l+1)}{\rho^2}
        +\frac{\lambda}{\rho}
    \right)R
    =
    0
    },
\]
где штрихи теперь обозначают производные по $\rho$.

При $\rho\to0$ регулярное решение ведёт себя как $\rho^l$, поэтому
выделим этот множитель:
\[
    R(\rho)=\rho^lP(\rho).
\]
Тогда
\[
    R'
    =
    l\rho^{l-1}P+\rho^lP',
\]
\[
    R''
    =
    l(l-1)\rho^{l-2}P
    +
    2l\rho^{l-1}P'
    +
    \rho^lP'',
\]
и после подстановки получаем
\[
    \boxed{
    P''
    +
    \frac{2(l+1)}{\rho}P'
    +
    \left(
        \frac{\lambda}{\rho}
        -
        \frac14
    \right)P
    =
    0
    }.
\]

\subsubsection{Асимптотика при больших \texorpdfstring{$\rho$}{rho}}

Для определения поведения $P$ при $\rho\to\infty$ введём
логарифмическую производную
\[
    Y=\frac{P'}{P}.
\]
Тогда
\[
    \frac{P''}{P}=Y'+Y^2,
\]
и уравнение принимает вид
\[
    Y'
    +
    Y^2
    +
    \frac{2(l+1)}{\rho}Y
    +
    \frac{\lambda}{\rho}
    -
    \frac14
    =
    0.
\]
При $\rho\to\infty$ остаётся
\[
    Y^2-\frac14\simeq0,
\]
поэтому
\[
    Y\simeq\pm\frac12.
\]
Следовательно,
\[
    P(\rho)\sim e^{\pm\rho/2}.
\]
Решение $e^{+\rho/2}$ не является нормируемым, поэтому физическое
решение должно содержать затухающий множитель:
\[
    \boxed{
        P(\rho)=e^{-\rho/2}L(\rho)
    }.
\]

Производные равны
\[
    P'
    =
    e^{-\rho/2}
    \left(
        L'-\frac12L
    \right),
\]
\[
    P''
    =
    e^{-\rho/2}
    \left(
        L''-L'+\frac14L
    \right).
\]
После подстановки получаем
\[
    \boxed{
        \rho L''
        +
        \bigl(2(l+1)-\rho\bigr)L'
        +
        (\lambda-l-1)L
        =
        0
    }.
\]

Это уравнение удобно записать в форме
\[
    \rho L''+(\beta-\rho)L'-\alpha L=0,
\]
где
\[
    \boxed{
        \beta=2(l+1),
        \qquad
        \alpha=l+1-\lambda
    }.
\]

\subsubsection{Разложение в ряд}

Ищем решение в виде степенного ряда
\[
    L(\rho)
    =
    \sum_{k=0}^{\infty}a_k\rho^k.
\]
Тогда
\[
    L'
    =
    \sum_{k=1}^{\infty}
        ka_k\rho^{k-1},
    \qquad
    L''
    =
    \sum_{k=2}^{\infty}
        k(k-1)a_k\rho^{k-2}.
\]
Подстановка в уравнение даёт рекуррентное соотношение
\[
    \boxed{
        a_{k+1}(k+1)(\beta+k)
        =
        a_k(\alpha+k)
    }.
\]
То есть
\[
    a_{k+1}
    =
    a_k
    \frac{\alpha+k}{(k+1)(\beta+k)}.
\]

Вводя возрастающий факториал
\[
    \alpha^{(k)}
    =
    \alpha(\alpha+1)\cdots(\alpha+k-1)
    =
    \frac{\Gamma(\alpha+k)}{\Gamma(\alpha)},
\]
аналогично для $\beta^{(k)}$, получаем
\[
    a_k
    =
    a_0
    \frac{\alpha^{(k)}}{\beta^{(k)}}
    \frac{1}{k!}.
\]
Поэтому решение представляется вырожденной гипергеометрической
функцией
\[
    \boxed{
    {}_1F_1(\alpha;\beta;\rho)
    =
    1+
    \sum_{k=1}^{\infty}
    \frac{\alpha^{(k)}}{\beta^{(k)}}
    \frac{\rho^k}{k!}
    }.
\]

Для связанного состояния степенной ряд должен оборваться. Это
происходит, если
\[
    \alpha=-n_r,
    \qquad
    n_r=0,1,2,\ldots
\]
Отсюда
\[
    l+1-\lambda=-n_r,
\]
то есть
\[
    \boxed{
        \lambda=n_r+l+1
    }.
\]
Величина
\[
    n=\lambda
\]
является главным квантовым числом:
\[
    \boxed{
        n=n_r+l+1=1,2,3,\ldots
    }.
\]

Подставляя $\lambda=n$ в выражение для энергии, получаем
\[
    E_n
    =
    -\frac{Ze^2}{2an^2}.
\]
Так как
\[
    a=\frac{\hbar^2}{\mu e^2Z},
\]
энергетический спектр можно записать как
\[
    \boxed{
        E_n
        =
        -\frac{\mu Z^2e^4}{2\hbar^2n^2}
        =
        -\frac{Z^2\,\mathrm{Ry}}{n^2}
    },
\]
где
\[
    \mathrm{Ry}
    =
    \frac{\mu e^4}{2\hbar^2}.
\]

% =========================================================
\section{Лекции 11--12}
% =========================================================

\subsection{Водородоподобный ион}

Продолжим решение радиального уравнения для связанных состояний
водородоподобного иона. После выделения асимптотического поведения
радиальной функции было получено уравнение
\[
    \rho L''+(2l+2-\rho)L'+(\lambda-l-1)L=0.
\]
Введём обозначения
\[
    \beta=2l+2,
    \qquad
    \alpha=l+1-\lambda.
\]
Тогда уравнение принимает стандартный вид
\[
    \boxed{
        \rho L''+(\beta-\rho)L'-\alpha L=0
    }.
\]
Радиальная функция при этом записывается как
\[
    \boxed{
        R(r)=C_{nl}e^{-\rho/2}\rho^lL(\rho)
    }.
\]

Решением уравнения для $L$ является вырожденная гипергеометрическая
функция
\[
    L(\rho)
    =
    {}_1F_1(\alpha;\beta;\rho)
    =
    1+
    \sum_{k=1}^{\infty}
    \frac{\alpha^{(k)}}{\beta^{(k)}}
    \frac{\rho^k}{k!},
\]
где используется обозначение
\[
    \alpha^{(k)}
    =
    \alpha(\alpha+1)\cdots(\alpha+k-1)
    =
    \frac{\Gamma(\alpha+k)}{\Gamma(\alpha)}.
\]

Чтобы исследовать поведение этого ряда при больших $\rho$, удобно
воспользоваться асимптотикой гамма-функции
\[
    \Gamma(x)
    \underset{x\to\infty}{\sim}
    \sqrt{\frac{2\pi}{x}}
    \left(\frac{x}{e}\right)^x.
\]
Для коэффициентов ряда
\[
    \frac{\alpha^{(k)}}{\beta^{(k)}k!}
    =
    \frac{\Gamma(\beta)}{\Gamma(\alpha)}
    \frac{\Gamma(\alpha+k)}
         {\Gamma(\beta+k)\Gamma(k+1)},
\]
и при $k\to\infty$ это приводит к асимптотическому поведению
\[
    {}_1F_1(\alpha;\beta;\rho)
    \sim
    \frac{\Gamma(\beta)}{\Gamma(\alpha)}
    \rho^{\alpha-\beta}e^\rho.
\]
Если ряд не обрывается, то множитель $e^\rho$ в $L(\rho)$ вместе с
$e^{-\rho/2}$ в $R(r)$ приводит к экспоненциальному росту радиальной
функции. Такое решение не является нормируемым.

Поэтому для связанных состояний необходимо, чтобы ряд обрывался:
\[
    \boxed{
        \alpha=-n_r,
        \qquad
        n_r=0,1,2,\ldots
    }.
\]
Здесь $n_r$ --- радиальное квантовое число. Оно определяет число узлов
радиальной волновой функции.

При $\alpha=-n_r$ имеем
\[
    \alpha^{(k)}
    =
    (-n_r)(-n_r+1)\cdots(-n_r+k-1),
\]
поэтому
\[
    \alpha^{(k)}=0,
    \qquad
    k\geq n_r+1,
\]
и гипергеометрический ряд превращается в полином степени $n_r$:
\[
\begin{aligned}
    L(\rho)
    &=
    {}_1F_1(-n_r;\beta;\rho) \\
    &=
    1+
    \sum_{k=1}^{n_r}
    \frac{(-n_r)^{(k)}}{\beta^{(k)}}
    \frac{\rho^k}{k!} \\
    &=
    1+
    \sum_{k=1}^{n_r}
    (-1)^k
    \frac{n_r!}{(n_r-k)!}
    \frac{\rho^k}{k!\,\beta^{(k)}}.
\end{aligned}
\]

Поскольку
\[
    \alpha=l+1-\lambda=-n_r,
\]
получаем
\[
    \boxed{
        \lambda=n_r+l+1\equiv n
    }.
\]
Таким образом, возникает главное квантовое число
\[
    \boxed{
        n=1,2,3,\ldots
    },
\]
а допустимые значения орбитального квантового числа при фиксированном
$n$ равны
\[
    \boxed{
        l=0,1,\ldots,n-1
    }.
\]

Связь между безразмерной переменной $\rho$ и радиусом имеет вид
\[
    r=\frac{a\lambda}{2}\rho
    =
    \frac{an}{2}\rho.
\]
Энергия равна
\[
    E
    =
    -\frac{Ze^2}{2a}\frac{1}{\lambda^2},
\]
поэтому после квантования $\lambda=n$
\[
    \boxed{
        E_n
        =
        -\frac{Ze^2}{2a}\frac{1}{n^2}
        =
        -\frac{Z^2\mathrm{Ry}}{n^2}
    }.
\]

Нормировочный коэффициент для радиальной функции можно записать как
\[
    \boxed{
    C_{nl}
    =
    \frac{2}
         {n^2a^{3/2}(2l+1)!}
    \sqrt{
        \frac{(n+l)!}{(n-l-1)!}
    }
    }.
\]

\subsubsection{Вырождение энергетических уровней}

При фиксированном $n$ возможны значения
\[
    l=0,1,\ldots,n-1.
\]
Для каждого $l$ магнитное квантовое число принимает значения
\[
    m=-l,-l+1,\ldots,l,
\]
то есть одному значению $l$ соответствует $2l+1$ состояний.

Поэтому полная кратность вырождения уровня с главным квантовым числом
$n$ равна
\[
    \boxed{
    g(n)
    =
    \sum_{l=0}^{n-1}(2l+1)
    =
    n^2
    }.
\]
Например,
\[
    n=1:
    \qquad
    l=0
    \qquad
    \Rightarrow
    \qquad
    1s,
\]
а для
\[
    n=2:
    \qquad
    l=0,1
    \qquad
    \Rightarrow
    \qquad
    2s,\;2p.
\]
То, что энергия кулоновской задачи зависит только от $n$, а не отдельно
от $l$, приводит к дополнительному вырождению уровней.

\subsubsection{Некоторые радиальные функции}

Для основного состояния $n=1$, $l=0$:
\[
    \boxed{
        R_{10}(r)
        =
        \frac{2}{a^{3/2}}e^{-r/a}
    }.
\]
Так как
\[
    Y_{00}=\frac{1}{\sqrt{4\pi}},
\]
полная волновая функция основного состояния равна
\[
    \boxed{
        \psi_{100}(\vec r)
        =
        \frac{1}{\sqrt{\pi a^3}}e^{-r/a}
    }.
\]

Для состояний с $n=2$:
\[
    \boxed{
    R_{21}(r)
    =
    \frac{1}{2\sqrt6\,a^{3/2}}
    \frac{r}{a}
    e^{-r/(2a)}
    },
\]
а для $l=0$
\[
    \boxed{
    R_{20}(r)
    =
    \frac{1}{\sqrt2\,a^{3/2}}
    e^{-r/(2a)}
    \left(
        1-\frac{r}{2a}
    \right)
    }.
\]
Для $R_{20}$ радиальное квантовое число равно
\[
    n_r=n-l-1=1,
\]
поэтому эта функция имеет один радиальный узел.

\subsubsection{Связь с присоединёнными полиномами Лагерра}

Полиномы, возникающие в радиальной части, можно записывать через
присоединённые полиномы Лагерра:
\[
    \boxed{
    L_n^{(\beta)}(x)
    =
    \frac{x^{-\beta}e^x}{n!}
    \left(
        \frac{d}{dx}
    \right)^n
    \left(
        e^{-x}x^{n+\beta}
    \right)
    }.
\]
Они связаны с вырожденной гипергеометрической функцией соотношением
\[
    \boxed{
    L_n^{(\beta)}(x)
    =
    \binom{n+\beta}{n}
    {}_1F_1(-n;\beta+1;x)
    }.
\]
Поэтому радиальную функцию водородоподобного иона можно выражать через
полином
\[
    L_{n-l-1}^{(2l+1)}.
\]
Здесь снова появляется комбинация
\[
    n-l-1=n_r,
\]
то есть степень соответствующего полинома совпадает с радиальным
квантовым числом.


% =========================================================
\subsubsection{Состояния при \texorpdfstring{$E>0$}{E>0}}
% =========================================================

При положительной энергии имеем непрерывный спектр. Энергию удобно
параметризовать волновым числом $k$:
\[
    \boxed{
        E=\frac{\hbar^2k^2}{2\mu}>0
    }.
\]
Формально сравнивая это выражение с формулой
\[
    E
    =
    -\frac{Ze^2}{2a}\frac{1}{n^2},
\]
получаем
\[
    n^2=-\frac{1}{k^2a^2},
    \qquad
    n\longrightarrow\frac{i}{ka}.
\]
Соответственно,
\[
    \rho
    =
    \frac{2r}{na}
    \longrightarrow
    -2ikr.
\]

Решение радиального уравнения можно записать в форме
\[
    \boxed{
    R_{kl}(r)
    =
    C_{kl}
    e^{ikr}
    (-2ikr)^l\,
    {}_1F_1
    \left(
        l+1-\frac{i}{ka};
        2l+2;
        -2ikr
    \right)
    }.
\]

При нормировке состояний непрерывного спектра на $\delta$-функцию
удобно использовать эквивалентную форму
\[
\boxed{
\begin{aligned}
    R_{kl}(r)
    &=
    \sqrt{\frac{2}{\pi}}\,
    \frac{k}{(2l+1)!}
    \left|
        \Gamma
        \left(
            l+1-\frac{i}{ka}
        \right)
    \right|
    e^{\frac{\pi}{2ka}}
    (2kr)^l e^{-ikr}
    \\
    &\qquad\times
    {}_1F_1
    \left(
        l+1+\frac{i}{ka};
        2l+2;
        2ikr
    \right).
\end{aligned}
}
\]

Для исследования поведения при больших $r$ снова удобно перейти к
функции
\[
    \chi(r)=rR(r).
\]
Она удовлетворяет уравнению
\[
    -\frac{\hbar^2}{2\mu}\chi''
    +
    \left(
        \frac{\hbar^2l(l+1)}{2\mu r^2}
        -
        \frac{Ze^2}{r}
    \right)\chi
    =
    \frac{\hbar^2k^2}{2\mu}\chi.
\]

Введём логарифмическую производную
\[
    Y=\frac{\chi'}{\chi}.
\]
Тогда
\[
    \frac{\chi''}{\chi}=Y'+Y^2,
\]
и радиальное уравнение принимает форму
\[
    \boxed{
    Y'+Y^2
    +
    \left(
        k^2
        -
        \frac{l(l+1)}{r^2}
        +
        \frac{2}{ar}
    \right)
    =
    0
    }.
\]

Полагая
\[
    \chi=Ae^{i\varphi},
    \qquad
    Y=i\varphi',
\]
при больших $r$ получаем
\[
    Y
    \simeq
    \pm i
    \left(
        k+\frac{1}{kar}
    \right).
\]
Следовательно,
\[
    \varphi(r)
    \simeq
    kr+\frac{1}{ak}\ln(kr)+\ldots
\]
и асимптотика радиальной функции имеет вид
\[
    \boxed{
    R_{kl}(r)
    \sim
    \frac{1}{r}
    \sin
    \left(
        kr
        +
        \frac{1}{ak}\ln(kr)
        -
        \frac{\pi l}{2}
        +
        \delta_l(k)
    \right)
    }.
\]
В отличие от короткодействующего потенциала, кулоновский потенциал
создаёт дополнительный логарифмический вклад в фазу.

Для фиксированного $l$ полнота радиальных связанных и непрерывных
состояний записывается как
\[
    \boxed{
    \frac{\delta(r-r')}{r^2}
    =
    \sum_n
        R_{nl}^*(r')R_{nl}(r)
    +
    \int_0^\infty
        dk\,
        R_{kl}^*(r')R_{kl}(r)
    }.
\]
Первая часть соответствует дискретным связанным состояниям, а интеграл
--- состояниям непрерывного спектра.


% =========================================================
\subsection{Абстрактная формулировка квантовой механики. Обозначения Дирака}
% =========================================================

В координатном представлении состояние частицы описывалось функцией
\[
    \psi(\vec r)\in L_2(\mathbb R^3).
\]
Абстрактно то же состояние обозначается кет-вектором
\[
    \boxed{
        |\psi\rangle\in\mathcal H
    },
\]
где $\mathcal H$ --- гильбертово пространство состояний.

Если
\[
    |\psi\rangle,\;|\varphi\rangle\in\mathcal H,
\]
то для любых комплексных чисел $\lambda,\mu$
\[
    \boxed{
        \lambda|\psi\rangle+\mu|\varphi\rangle
        \in\mathcal H
    }.
\]
Таким образом, пространство состояний является линейным.

\subsubsection{Скалярное произведение и норма}

В $\mathcal H$ определено скалярное произведение
\[
    \langle\varphi|\psi\rangle.
\]
Оно обладает свойствами
\[
    \langle\psi|\psi\rangle\geq0,
\]
причём
\[
    \langle\psi|\psi\rangle=0
    \quad\Longleftrightarrow\quad
    |\psi\rangle=0,
\]
а также
\[
    \boxed{
        \langle\varphi|\psi\rangle
        =
        \langle\psi|\varphi\rangle^*
    }.
\]
При принятом соглашении скалярное произведение линейно по второму
аргументу:
\[
\boxed{
\langle
\varphi
|
\lambda_1\psi_1+\lambda_2\psi_2
\rangle
=
\lambda_1\langle\varphi|\psi_1\rangle
+
\lambda_2\langle\varphi|\psi_2\rangle
}.
\]

Норма состояния определяется формулой
\[
    \boxed{
        \|\psi\|^2
        =
        \langle\psi|\psi\rangle
    }.
\]
Кет-вектору $|\psi\rangle\in\mathcal H$ соответствует бра-вектор
\[
    \langle\psi|\in\mathcal H^*,
\]
то есть элемент сопряжённого пространства.

\subsubsection{Ортонормированный базис}

Пусть в $\mathcal H$ выбран счётный ортонормированный базис
\[
    \{|e_n\rangle\}_{n=1}^{\infty},
\]
для которого
\[
    \boxed{
        \langle e_n|e_m\rangle
        =
        \delta_{nm}
    }.
\]
Тогда произвольное состояние раскладывается по базису:
\[
    \boxed{
        |\psi\rangle
        =
        \sum_{n=1}^{\infty}
        c_n|e_n\rangle
    },
\]
где
\[
    \boxed{
        c_n=\langle e_n|\psi\rangle
    }.
\]
Норма при этом равна
\[
    \boxed{
        \langle\psi|\psi\rangle
        =
        \sum_{n=1}^{\infty}|c_n|^2
        <\infty
    }.
\]

Если состояние не нормировано, вероятность получить при измерении
состояние, соответствующее базисному вектору $|e_n\rangle$, равна
\[
    \boxed{
        P_n
        =
        \frac{|c_n|^2}{\|\psi\|^2}
    }.
\]
Для нормированного состояния
\[
    \sum_nP_n=1.
\]

Физическому квантовому состоянию соответствует не конкретный вектор,
а луч в гильбертовом пространстве:
\[
    \boxed{
        z|\psi\rangle,
        \qquad
        z\in\mathbb C\setminus\{0\}
    }.
\]
То есть умножение всего вектора состояния на ненулевой комплексный
коэффициент не меняет физического состояния.

\subsubsection{Представление через пространство \texorpdfstring{$l_2$}{l2}}

Коэффициентам разложения
\[
    |\psi\rangle
    =
    \sum_{n=1}^{\infty}c_n|e_n\rangle
\]
можно сопоставить последовательность
\[
    \begin{pmatrix}
        c_1\\
        c_2\\
        \vdots\\
        c_n\\
        \vdots
    \end{pmatrix}.
\]
Пространство таких последовательностей обозначается $l_2$:
\[
    \boxed{
    l_2
    =
    \left\{
        \{c_n\}_{n=1}^{\infty}
        \;:\;
        \sum_{n=1}^{\infty}|c_n|^2<\infty
    \right\}
    }.
\]

Если
\[
    |\psi\rangle\longleftrightarrow\{c_n\},
    \qquad
    |\varphi\rangle\longleftrightarrow\{d_n\},
\]
то
\[
    \boxed{
        \langle\varphi|\psi\rangle
        =
        \sum_{n=1}^{\infty}d_n^*c_n
    }.
\]
Таким образом, выбор ортонормированного базиса позволяет представить
состояния сепарабельного гильбертова пространства последовательностями
из $l_2$.

\subsubsection{Тензорное произведение пространств}

Если имеются две квантовые системы с пространствами состояний
$\mathcal H_A$ и $\mathcal H_B$, то пространство состояний составной
системы строится как
\[
    \boxed{
        \mathcal H
        =
        \mathcal H_A\otimes\mathcal H_B
    }.
\]
Если
\[
    \{|e_i^A\rangle\}
    \quad\text{и}\quad
    \{|e_j^B\rangle\}
\]
--- базисы отдельных пространств, то базис тензорного произведения
образуют векторы
\[
    \boxed{
        |e_i^A\rangle\otimes|e_j^B\rangle
    }.
\]
Само тензорное произведение пространств не зависит от конкретного
выбора базиса.

Например,
\[
    \boxed{
        L_2(\mathbb R)\otimes L_2(\mathbb R)
        \simeq
        L_2(\mathbb R^2)
    }.
\]
Функции $\psi(x)$ и $\varphi(y)$ относятся к отдельным пространствам,
а состояние составной системы в общем случае описывается функцией
\[
    \chi(x,y).
\]


% =========================================================
\subsubsection{Линейные операторы}
% =========================================================

Линейный оператор
\[
    \hat A:\mathcal H\to\mathcal H
\]
должен удовлетворять условию
\[
\boxed{
\hat A
\left(
    \lambda_1|\psi_1\rangle+
    \lambda_2|\psi_2\rangle
\right)
=
\lambda_1\hat A|\psi_1\rangle+
\lambda_2\hat A|\psi_2\rangle
}.
\]

Пусть
\[
    |\psi\rangle
    =
    \sum_{m=1}^{\infty}c_m|e_m\rangle,
\]
и
\[
    |\varphi\rangle
    =
    \hat A|\psi\rangle.
\]
Тогда
\[
    |\varphi\rangle
    =
    \sum_{m=1}^{\infty}
        c_m\hat A|e_m\rangle.
\]
Разложим результат по тому же базису:
\[
    |\varphi\rangle
    =
    \sum_{n=1}^{\infty}
        d_n|e_n\rangle.
\]
Коэффициенты равны
\[
\begin{aligned}
    d_n
    &=
    \langle e_n|\varphi\rangle
    =
    \langle e_n|\hat A|\psi\rangle \\
    &=
    \sum_{m=1}^{\infty}
    \langle e_n|\hat A|e_m\rangle c_m.
\end{aligned}
\]
Вводя
\[
    \boxed{
        A_{nm}
        =
        \langle e_n|\hat A|e_m\rangle
    },
\]
получаем
\[
    \boxed{
        d_n
        =
        \sum_{m=1}^{\infty}A_{nm}c_m
    }.
\]
Таким образом, числа $A_{nm}$ образуют матричное представление
оператора $\hat A$ в выбранном базисе.

Величины вида
\[
    \langle\varphi|\hat A|\psi\rangle
\]
называются матричными элементами оператора.

\subsubsection{Непрерывные и ограниченные операторы}

Линейный оператор $\hat A$ называется непрерывным, если из
\[
    |x_n\rangle\longrightarrow|x\rangle
\]
следует
\[
    \boxed{
        \hat A|x_n\rangle
        \longrightarrow
        \hat A|x\rangle
    }.
\]

Для линейного оператора непрерывность эквивалентна ограниченности:
существует постоянная $C$, такая что
\[
    \boxed{
        \|\hat A|x\rangle\|
        \leq
        C\|\,|x\rangle\|
    }.
\]
Если линейный оператор непрерывен в одной точке, то он непрерывен на
всём пространстве.

\subsubsection{Неограниченные операторы и область определения}

Неограниченный оператор, вообще говоря, не может быть определён на
всём гильбертовом пространстве. Поэтому оператор необходимо задавать
вместе с его областью определения:
\[
    \boxed{
        \hat A
        =
        \{A,\mathcal D(\hat A)\}
    },
    \qquad
    \mathcal D(\hat A)\subset\mathcal H.
\]
Оператор действует как
\[
    \hat A:
    \mathcal D(\hat A)\to\mathcal H.
\]
Обычно рассматриваются плотно определённые операторы:
\[
    \boxed{
        \overline{\mathcal D(\hat A)}
        =
        \mathcal H
    }.
\]

Для неограниченных операторов область определения является частью
самого определения оператора. Поэтому два оператора считаются равными,
если совпадают не только их действия, но и области определения.

Оператор $\hat A$ называется замкнутым, если из
\[
    |x_n\rangle\in\mathcal D(\hat A),
    \qquad
    |x_n\rangle\to|x\rangle,
\]
и
\[
    \hat A|x_n\rangle\to|y\rangle
\]
следует
\[
    \boxed{
        |x\rangle\in\mathcal D(\hat A),
        \qquad
        \hat A|x\rangle=|y\rangle
    }.
\]

Для двух операторов пишут
\[
    \hat A\subseteq\hat B,
\]
если
\[
    \mathcal D(\hat A)\subseteq\mathcal D(\hat B)
\]
и
\[
    \boxed{
        \hat B|x\rangle
        =
        \hat A|x\rangle,
        \qquad
        |x\rangle\in\mathcal D(\hat A)
    }.
\]
Равенство операторов требует
\[
    \boxed{
        \mathcal D(\hat A)=\mathcal D(\hat B)
    }
\]
и совпадения их действия на этой области.

Для произведения операторов $\hat A\hat B$ необходимо, чтобы после
действия $\hat B$ полученный вектор принадлежал области определения
$\hat A$:
\[
    \mathcal D(\hat A\hat B)
    =
    \left\{
        |x\rangle\in\mathcal D(\hat B):
        \hat B|x\rangle\in\mathcal D(\hat A)
    \right\}.
\]

Замкнутый линейный оператор, определённый на всём гильбертовом
пространстве, является ограниченным.


% =========================================================
\subsubsection{Спектр оператора}
% =========================================================

Собственный вектор оператора удовлетворяет уравнению
\[
    \boxed{
        \hat A|\psi_i\rangle
        =
        \lambda_i|\psi_i\rangle
    },
\]
или
\[
    \boxed{
        (\hat A-\lambda_i\hat I)
        |\psi_i\rangle
        =
        0
    }.
\]

Для описания спектра рассматривается оператор
\[
    \hat A-\lambda\hat I.
\]
Если существует обратный оператор
\[
    \boxed{
        \hat R_A(\lambda)
        =
        (\hat A-\lambda\hat I)^{-1}
    }
\]
и он является ограниченным оператором на $\mathcal H$, то
$\lambda$ принадлежит резольвентному множеству:
\[
    \lambda\in\rho(\hat A).
\]
Спектр оператора определяется как дополнение резольвентного множества:
\[
    \boxed{
        \sigma(\hat A)
        =
        \mathbb C\setminus\rho(\hat A)
    }.
\]

Все собственные значения оператора принадлежат его спектру.
Совокупность собственных значений образует точечный спектр.

Изолированные точки спектра образуют дискретную часть спектра.

Если
\[
    \overline{
        \operatorname{Ran}
        (\hat A-\lambda\hat I)
    }
    \neq
    \mathcal H,
\]
то соответствующие точки относятся к остаточному спектру. В
квантово-механических задачах такой спектр практически не возникает.

Если же
\[
    \boxed{
    \overline{
        \operatorname{Ran}
        (\hat A-\lambda\hat I)
    }
    =
    \mathcal H
    },
\]
но $\lambda$ не принадлежит резольвентному множеству и не является
собственным значением, возникает непрерывный спектр.

Здесь
\[
    \operatorname{Im}
    =
    \operatorname{Ran}
\]
обозначает область значений оператора.


% =========================================================
\subsection{Эрмитово сопряжение операторов}
% =========================================================

Для оператора $\hat A$ вводится эрмитово сопряжённый оператор
$\hat A^\dagger$. В случае ограниченных операторов он определяется
условием
\[
    \boxed{
    \langle\varphi|\hat A|\psi\rangle
    =
    \langle\psi|\hat A^\dagger|\varphi\rangle^*
    }
\]
для любых
\[
    |\varphi\rangle,\;|\psi\rangle\in\mathcal H.
\]

В матричном представлении
\[
    \boxed{
        (A^\dagger)_{nm}
        =
        A_{mn}^*
    }.
\]
Таким образом, матрица сопряжённого оператора получается комплексным
сопряжением и транспонированием матрицы исходного оператора.

Если
\[
    \langle\varphi|\hat A|\psi\rangle
    =
    \langle\psi|\hat A|\varphi\rangle^*,
\]
то оператор обладает свойством эрмитовости. Для неограниченных
операторов при этом необходимо отдельно учитывать области определения.

\subsubsection{Сопряжённый оператор для неограниченного оператора}

Пусть
\[
    |\psi\rangle\in\mathcal D(\hat A).
\]
Для фиксированного $|\varphi\rangle$ выражение
\[
    \langle\varphi|\hat A|\psi\rangle
\]
можно рассматривать как линейный функционал от $|\psi\rangle$:
\[
    \eta_\varphi:
    \mathcal D(\hat A)\to\mathbb C.
\]

Область определения сопряжённого оператора состоит из тех векторов
$|\varphi\rangle\in\mathcal H$, для которых существует вектор
$|\eta\rangle\in\mathcal H$, удовлетворяющий
\[
    \boxed{
    \langle\varphi|\hat A|\psi\rangle
    =
    \langle\eta|\psi\rangle
    \qquad
    \forall\,|\psi\rangle\in\mathcal D(\hat A)
    }.
\]
Иными словами,
\[
\boxed{
\mathcal D(\hat A^\dagger)
=
\left\{
|\varphi\rangle\in\mathcal H:
\exists\,|\eta\rangle\in\mathcal H
\text{ такое, что }
\langle\varphi|\hat A|\psi\rangle
=
\langle\eta|\psi\rangle
\;
\forall|\psi\rangle\in\mathcal D(\hat A)
\right\}
}.
\]

После этого действие сопряжённого оператора определяется как
\[
    \boxed{
        \hat A^\dagger|\varphi\rangle
        =
        |\eta\rangle
    }.
\]
Тогда
\[
    \boxed{
    \langle\varphi|\hat A|\psi\rangle
    =
    \langle\psi|\hat A^\dagger|\varphi\rangle^*
    }.
\]

Оператор $\hat A$ называется эрмитовым, если
\[
    \boxed{
        \hat A\subseteq\hat A^\dagger
    }.
\]
Это означает, что
\[
    \mathcal D(\hat A)
    \subseteq
    \mathcal D(\hat A^\dagger)
\]
и для любых
\[
    |\varphi\rangle,\;|\psi\rangle\in\mathcal D(\hat A)
\]
выполняется
\[
    \boxed{
        \langle\varphi|\hat A|\psi\rangle
        =
        \langle\psi|\hat A|\varphi\rangle^*
    }.
\]

Если же
\[
    \boxed{
        \hat A=\hat A^\dagger
    },
\]
то есть совпадают как действия операторов, так и их области
определения,
\[
    \boxed{
        \mathcal D(\hat A)
        =
        \mathcal D(\hat A^\dagger)
    },
\]
оператор называется \emph{самосопряжённым}.

% =========================================================
\section{Лекции 13--14}
% =========================================================

\subsection{Самосопряжённые и симметричные операторы}

Пусть
\[
    \hat A:\mathcal D(\hat A)\to\mathcal H.
\]
Самосопряжённый оператор удовлетворяет
\[
    \boxed{
        \hat A^\dagger=\hat A
    },
\]
то есть совпадают как действия операторов, так и их области
определения. Для симметричного (эрмитова) оператора выполняется более
слабое условие
\[
    \boxed{
        \hat A\subseteq\hat A^\dagger
    }.
\]

Если область определения $\hat A$ плотна в $\mathcal H$,
\[
    \overline{\mathcal D(\hat A)}=\mathcal H,
\]
то сопряжённый оператор $\hat A^\dagger$ является замкнутым.

Действительно, пусть
\[
    |\psi_n\rangle\in\mathcal D(\hat A^\dagger),
    \qquad
    |\psi_n\rangle\longrightarrow|\psi\rangle,
\]
и одновременно
\[
    \hat A^\dagger|\psi_n\rangle
    \longrightarrow
    |\eta\rangle.
\]
Для любого $|\varphi\rangle\in\mathcal D(\hat A)$ имеем
\[
    \langle\varphi|\psi_n\rangle
    \longrightarrow
    \langle\varphi|\psi\rangle.
\]
По определению сопряжённого оператора
\[
    \langle\varphi|\hat A^\dagger|\psi_n\rangle
    =
    \langle\psi_n|\hat A|\varphi\rangle^*.
\]
При переходе к пределу
\[
    \langle\varphi|\eta\rangle
    =
    \langle\psi|\hat A|\varphi\rangle^*,
    \qquad
    \forall\,|\varphi\rangle\in\mathcal D(\hat A).
\]
Следовательно,
\[
    |\psi\rangle\in\mathcal D(\hat A^\dagger),
    \qquad
    \hat A^\dagger|\psi\rangle=|\eta\rangle,
\]
что и означает замкнутость $\hat A^\dagger$.

Если оператор $\hat A$ инъективен, то существует обратный оператор
$\hat A^{-1}$, причём
\[
    \boxed{
        \mathcal D(\hat A^{-1})
        =
        \operatorname{Im}\hat A
    }.
\]
Если $\hat A$ замкнут, то и $\hat A^{-1}$ является замкнутым на своей
области определения.

Если
\[
    \lambda\in\rho(\hat A),
\]
то существует ограниченный оператор
\[
    (\hat A-\lambda\hat I)^{-1}.
\]
Отсюда следует, что $\hat A-\lambda\hat I$, а значит и $\hat A$,
является замкнутым. В частности, самосопряжённые операторы замкнуты.


% =========================================================
\subsubsection{Вещественность спектра самосопряжённого оператора}
% =========================================================

Спектр самосопряжённого оператора вещественен:
\[
    \boxed{
        \sigma(\hat A)\subset\mathbb R
    }.
\]

Для собственного значения это видно непосредственно. Если
\[
    \hat A|\psi\rangle
    =
    \lambda|\psi\rangle,
\]
то
\[
    \langle\psi|\hat A|\psi\rangle
    =
    \lambda\langle\psi|\psi\rangle.
\]
С другой стороны, поскольку $\hat A^\dagger=\hat A$,
\[
    \langle\psi|\hat A|\psi\rangle
    =
    \langle\psi|\hat A^\dagger|\psi\rangle
    =
    \langle\psi|\hat A|\psi\rangle^*.
\]
Поэтому $\lambda=\lambda^*$ и
\[
    \lambda\in\mathbb R.
\]

Покажем более общее утверждение для всего спектра. Пусть
\[
    \lambda=\alpha+i\mu,
    \qquad
    \alpha,\mu\in\mathbb R,
    \qquad
    \mu\neq0.
\]
Для любого $|\psi\rangle\in\mathcal D(\hat A)$ рассмотрим
\[
    \|(\hat A-\lambda\hat I)|\psi\rangle\|^2.
\]
Так как
\[
    \hat A-\lambda\hat I
    =
    \hat A-\alpha\hat I-i\mu\hat I,
\]
получаем
\begin{align*}
    \|(\hat A-\lambda\hat I)|\psi\rangle\|^2
    &=
    \langle\psi|
    (\hat A-\alpha\hat I+i\mu\hat I)
    (\hat A-\alpha\hat I-i\mu\hat I)
    |\psi\rangle \\
    &=
    \|(\hat A-\alpha\hat I)|\psi\rangle\|^2
    +
    |\mu|^2\|\psi\|^2.
\end{align*}
Следовательно,
\[
    \boxed{
    \|(\hat A-\lambda\hat I)|\psi\rangle\|^2
    \geq
    |\mu|^2\|\psi\|^2
    }.
\]

В частности,
\[
    (\hat A-\lambda\hat I)|\psi\rangle=0
    \quad\Longrightarrow\quad
    |\psi\rangle=0,
\]
поэтому $\hat A-\lambda\hat I$ инъективен.

Кроме того, из оценки следует
\[
    \|\psi\|
    \leq
    \frac{1}{|\mu|}
    \|(\hat A-\lambda\hat I)|\psi\rangle\|.
\]

Покажем, что
\[
    \operatorname{Im}(\hat A-\lambda\hat I)
\]
плотно в $\mathcal H$. Пусть некоторый $|\varphi\rangle$ ортогонален
всему образу:
\[
    \langle\varphi|
    (\hat A-\lambda\hat I)
    |\psi\rangle
    =
    0,
    \qquad
    \forall\,|\psi\rangle\in\mathcal D(\hat A).
\]
Тогда
\[
    |\varphi\rangle\in\mathcal D(\hat A^\dagger)
\]
и
\[
    (\hat A^\dagger-\lambda^*\hat I)|\varphi\rangle=0.
\]
Но $\hat A^\dagger=\hat A$, а $\operatorname{Im}\lambda^*=-\mu\neq0$,
поэтому
\[
    |\varphi\rangle=0.
\]
Следовательно,
\[
    \overline{
        \operatorname{Im}(\hat A-\lambda\hat I)
    }
    =
    \mathcal H.
\]

Остаётся убедиться, что образ замкнут. Пусть
\[
    (\hat A-\lambda\hat I)|\psi_n\rangle
    \longrightarrow
    |\varphi\rangle.
\]
Тогда
\[
    \|\psi_n-\psi_m\|
    \leq
    \frac{1}{|\mu|}
    \|
        (\hat A-\lambda\hat I)
        (|\psi_n\rangle-|\psi_m\rangle)
    \|,
\]
поэтому $\{|\psi_n\rangle\}$ является фундаментальной
последовательностью:
\[
    |\psi_n\rangle\longrightarrow|\psi\rangle.
\]
Так как $\hat A$ замкнут,
\[
    |\psi\rangle\in\mathcal D(\hat A),
    \qquad
    (\hat A-\lambda\hat I)|\psi\rangle
    =
    |\varphi\rangle.
\]
Значит,
\[
    \operatorname{Im}(\hat A-\lambda\hat I)=\mathcal H.
\]
Таким образом, при $\operatorname{Im}\lambda\neq0$ оператор
$\hat A-\lambda\hat I$ имеет ограниченный обратный оператор, то есть
\[
    \lambda\in\rho(\hat A).
\]
Следовательно,
\[
    \sigma(\hat A)\subset\mathbb R.
\]

Для симметричного оператора, который не является самосопряжённым,
возможны различные варианты поведения невещественной части спектра:
она может отсутствовать либо быть связана с верхней, нижней или обеими
полуплоскостями комплексной плоскости.


% =========================================================
\subsubsection{Некоторые свойства сопряжённых операторов}
% =========================================================

Если $\hat A$ --- замкнутый плотно определённый оператор, то
\[
    \boxed{
        \hat A^\dagger\hat A
        \text{ является самосопряжённым и неотрицательным}
    }.
\]

Для замкнутого плотно определённого оператора
\[
    \boxed{
        \hat A^{\dagger\dagger}=\hat A
    }.
\]

Для произведения операторов в общем случае выполняется включение
\[
    \boxed{
        (\hat A\hat B)^\dagger
        \supseteq
        \hat B^\dagger\hat A^\dagger
    }.
\]
Если $\hat A$ и $\hat B$ ограничены, то области определения не создают
дополнительных ограничений и
\[
    \boxed{
        (\hat A\hat B)^\dagger
        =
        \hat B^\dagger\hat A^\dagger
    }.
\]

Аналогично,
\[
    \boxed{
    (\lambda\hat A+\mu\hat B)^\dagger
    \supseteq
    \lambda^*\hat A^\dagger
    +
    \mu^*\hat B^\dagger
    }.
\]


% =========================================================
\subsubsection{Проекторы}
% =========================================================

Оператор $\hat P$ называется проектором, если
\[
    \boxed{
        \hat P^2=\hat P
    }.
\]
Если дополнительно
\[
    \boxed{
        \hat P^\dagger=\hat P
    },
\]
то $\hat P$ является ортогональным проектором.

Любой вектор можно разложить на две части:
\[
    |\psi\rangle
    =
    \hat P|\psi\rangle
    +
    (\hat I-\hat P)|\psi\rangle.
\]
Первая часть принадлежит подпространству, на которое проектирует
$\hat P$, а вторая --- его ортогональному дополнению.

Для нормированного вектора $|\psi\rangle$ проектор на одномерное
подпространство, натянутое на этот вектор, имеет вид
\[
    \boxed{
        \hat P_\psi
        =
        |\psi\rangle\langle\psi|
    }.
\]
Действительно,
\[
    \hat P_\psi|\varphi\rangle
    =
    |\psi\rangle
    \langle\psi|\varphi\rangle.
\]


% =========================================================
\subsubsection{Пример: оператор импульса на отрезке}
% =========================================================

Рассмотрим пространство
\[
    L_2([0,1])
\]
и дифференциальный оператор
\[
    \boxed{
        \hat p=-i\hbar\frac{d}{dx}
    }.
\]

Сначала выберем область определения
\[
    \mathcal D(\hat p_1)
    =
    \left\{
        \psi\in AC[0,1]:
        \psi(0)=\psi(1)=0
    \right\},
\]
где $AC[0,1]$ обозначает абсолютно непрерывные функции на отрезке.

Формальное уравнение на собственные значения
\[
    \hat p|\psi\rangle
    =
    \lambda|\psi\rangle
\]
даёт
\[
    \psi(x)
    =
    C e^{i\lambda x/\hbar}.
\]
Однако ненулевое решение такого вида не может одновременно
удовлетворять условиям
\[
    \psi(0)=\psi(1)=0.
\]
Поэтому с указанной областью определения оператор $\hat p_1$ не
является самосопряжённым.

Найдём сопряжённый оператор. Для
\[
    \psi\in\mathcal D(\hat p_1)
\]
интегрирование по частям даёт
\begin{align*}
    \langle\varphi|\hat p_1|\psi\rangle
    &=
    -i\hbar
    \int_0^1
        \varphi^*(x)\psi'(x)\,dx \\
    &=
    -i\hbar
    \left.
        \varphi^*(x)\psi(x)
    \right|_0^1
    +
    i\hbar
    \int_0^1
        \frac{d\varphi^*}{dx}\psi(x)\,dx.
\end{align*}
Из-за граничных условий на $\psi$ первый член исчезает. Поэтому
\[
    \mathcal D(\hat p_1^\dagger)
    =
    AC[0,1],
\]
то есть для сопряжённого оператора граничные условия отсутствуют, а
действие остаётся тем же:
\[
    \hat p_1^\dagger
    =
    -i\hbar\frac{d}{dx}.
\]
Таким образом,
\[
    \mathcal D(\hat p_1)
    \subset
    \mathcal D(\hat p_1^\dagger),
\]
и $\hat p_1$ является симметричным оператором.

Чтобы получить самосопряжённый оператор импульса, необходимо выбрать
подходящие граничные условия. Из интегрирования по частям следует, что
граничный член должен обращаться в нуль. В частности, достаточно
потребовать
\[
    |\psi(0)|^2=|\psi(1)|^2.
\]
Это приводит к семейству условий
\[
    \boxed{
        \psi(1)=e^{i\alpha}\psi(0),
        \qquad
        \alpha\in[0,2\pi)
    }.
\]

Определим оператор $\hat p_\alpha$ той же формулой
\[
    \hat p_\alpha
    =
    -i\hbar\frac{d}{dx},
\]
но на области
\[
    \boxed{
    \mathcal D(\hat p_\alpha)
    =
    \left\{
        \psi\in AC[0,1]:
        \psi(1)=e^{i\alpha}\psi(0)
    \right\}
    }.
\]
Для этой области оператор является самосопряжённым.

Уравнение
\[
    \hat p_\alpha|\psi\rangle
    =
    \lambda|\psi\rangle
\]
имеет решение
\[
    \psi(x)=Ce^{i\lambda x/\hbar}.
\]
Граничное условие требует
\[
    e^{i\lambda/\hbar}=e^{i\alpha},
\]
откуда
\[
    \boxed{
        \lambda
        =
        \hbar(2\pi n+\alpha),
        \qquad
        n\in\mathbb Z
    }.
\]
Все собственные значения вещественны, как и должно быть для
самосопряжённого оператора.

При этом
\[
    \hat p_\alpha
    \subset
    \hat p_1^\dagger.
\]

В качестве других примеров операторных выражений можно рассматривать
симметризованную комбинацию
\[
    \boxed{
        \hat x^3\hat p_x
        +
        \hat p_x\hat x^3
    }
\]
и радиальный оператор импульса, формально записываемый как
\[
    \boxed{
        \hat p_r
        =
        -i\hbar\frac{d}{dr}
    }.
\]


% =========================================================
\subsection{Спектральная теорема}
% =========================================================

Рассмотрим сначала самосопряжённый оператор $\hat A$ в
конечномерном гильбертовом пространстве:
\[
    \dim\mathcal H=N<\infty.
\]
Пусть
\[
    \{|e_n\rangle\}_{n=1}^{N}
\]
--- ортонормированный базис,
\[
    \langle e_n|e_m\rangle
    =
    \delta_{nm}.
\]

У самосопряжённого оператора существует ортонормированный базис из
собственных векторов:
\[
    \boxed{
        \hat A|e_n\rangle
        =
        \lambda_n|e_n\rangle,
        \qquad
        \lambda_n\in\mathbb R
    }.
\]
Поэтому
\[
    \boxed{
        \hat A
        =
        \sum_{n=1}^{N}
        \lambda_n
        |e_n\rangle\langle e_n|
    }.
\]
Действие оператора на произвольный вектор имеет вид
\[
    \hat A|\psi\rangle
    =
    \sum_{n=1}^{N}
        \lambda_n
        |e_n\rangle
        \langle e_n|\psi\rangle.
\]

Для единичного оператора
\[
    \boxed{
        \hat I
        =
        \sum_{n=1}^{N}
        |e_n\rangle\langle e_n|
    },
\]
поэтому
\[
    |\psi\rangle
    =
    \sum_{n=1}^{N}
        |e_n\rangle
        \langle e_n|\psi\rangle.
\]

Если некоторые собственные значения вырождены, удобно объединить
собственные векторы с одинаковым $\lambda_n$ в одно собственное
подпространство. Тогда
\[
    \boxed{
        \hat A
        =
        \sum_n
        \lambda_n\hat P_n
    },
\]
где теперь $\lambda_n$ --- различные собственные значения, а
$\hat P_n$ --- ортогональный проектор на соответствующее собственное
подпространство.

При этом
\[
    \boxed{
        \hat I=\sum_n\hat P_n
    },
\]
а если
\[
    \{|e_{nk}\rangle\}_k
\]
--- ортонормированный базис собственного подпространства с
собственным значением $\lambda_n$, то
\[
    \boxed{
        \hat P_n
        =
        \sum_k
        |e_{nk}\rangle\langle e_{nk}|
    }.
\]
Набор собственных значений образует спектр:
\[
    \{\lambda_n\}
    =
    \sigma(\hat A),
\]
и
\[
    \boxed{
        \hat A\hat P_n
        =
        \lambda_n\hat P_n
    }.
\]


% =========================================================
\subsubsection{Правило Борна и спектральное разложение}
% =========================================================

Пусть состояние нормировано:
\[
    \langle\psi|\psi\rangle=1.
\]
Вероятность получить при измерении собственное значение $\lambda_n$
равна квадрату нормы проекции состояния на соответствующее собственное
подпространство:
\[
    \boxed{
        P_n(\lambda_n)
        =
        \|\hat P_n|\psi\rangle\|^2
    }.
\]
Поскольку $\hat P_n^\dagger=\hat P_n$ и
$\hat P_n^2=\hat P_n$,
\begin{align*}
    P_n(\lambda_n)
    &=
    \langle\psi|
        \hat P_n^\dagger\hat P_n
    |\psi\rangle \\
    &=
    \boxed{
        \langle\psi|\hat P_n|\psi\rangle
    }.
\end{align*}

Среднее значение наблюдаемой $\hat A$ равно
\[
    \boxed{
    \langle\psi|\hat A|\psi\rangle
    =
    \sum_n
        \lambda_n
        P_n(\lambda_n)
    }.
\]

Для описания распределения результатов измерения введём функцию
распределения
\[
    M_\psi(\lambda)
    =
    P_\psi(A\leq\lambda).
\]
Она удовлетворяет условиям
\[
    \boxed{
        M_\psi(-\infty)=0,
        \qquad
        M_\psi(+\infty)=1
    }.
\]
Для дискретного спектра
\[
    \boxed{
    M_\psi(\lambda)
    =
    \sum_n
        \langle\psi|\hat P_n|\psi\rangle
        \Theta(\lambda-\lambda_n)
    },
\]
где $\Theta$ --- ступенчатая функция.

В дифференциальной форме
\[
    \boxed{
    dM_\psi(\lambda)
    =
    \sum_n
        \langle\psi|\hat P_n|\psi\rangle
        \delta(\lambda-\lambda_n)\,d\lambda
    }.
\]
Таким образом, дискретная сумма естественным образом переходит к
интегральной записи
\[
    \sum_n
    \quad\longrightarrow\quad
    \int dM_\psi(\lambda).
\]

В более общем случае $M_\psi$ рассматривается как вероятностная мера:
\[
    M_\psi:\{\Omega\}\to\mathbb R_+.
\]
Она удовлетворяет стандартным свойствам
\[
    M_\psi(\Omega)\geq0,
    \qquad
    M_\psi(\varnothing)=0,
    \qquad
    M_\psi(\mathbb R)=1,
\]
а для попарно непересекающихся множеств $\Omega_i$
\[
    \boxed{
    M_\psi
    \left(
        \bigcup_i\Omega_i
    \right)
    =
    \sum_i
        M_\psi(\Omega_i)
    }.
\]

Функция распределения связана с этой мерой как
\[
    M_\psi(\lambda)
    =
    M_\psi((-\infty,\lambda]).
\]

В конечномерном случае среднее значение записывалось как
\[
    \langle\psi|\hat A|\psi\rangle
    =
    \sum_n
        \lambda_n
        \langle\psi|\hat P_n|\psi\rangle.
\]
В спектральной форме эта запись обобщается до интеграла
\[
    \boxed{
    \langle\psi|\hat A|\psi\rangle
    =
    \int_{\mathbb R}
        \lambda\,dM_\psi(\lambda)
    }.
\]

% =========================================================
\subsubsection{Проекторнозначная мера}
% =========================================================

Для перехода от дискретного спектра к общему случаю вводится
\emph{проекторнозначная мера}. Каждому измеримому множеству
$\Omega\subset\mathbb R$ сопоставляется ортогональный проектор
\[
    \boxed{
        \hat P(\Omega):\{\Omega\}\longrightarrow
        \{\text{проекторы в }\mathcal H\}
    }.
\]
Для неё выполняются условия
\[
    \boxed{
        \hat P(\mathbb R)=\hat I,
        \qquad
        \hat P(\varnothing)=0
    }.
\]
Если множества $\Omega_i$ попарно не пересекаются, то
\[
    \boxed{
        \hat P\left(\bigcup_i\Omega_i\right)
        =
        \sum_i\hat P(\Omega_i)
    }.
\]
Сумма здесь понимается в соответствующем операторном смысле.

Особенно часто используется семейство проекторов
\[
    \boxed{
        \hat P_\lambda
        =
        \hat P((-\infty,\lambda])
    }.
\]
Оно соответствует всем значениям наблюдаемой, не превосходящим
$\lambda$.

Для любого состояния $|\psi\rangle\in\mathcal H$ проекторнозначная
мера порождает обычную вероятностную меру
\[
    \boxed{
        M_\psi(\Omega)
        =
        \langle\psi|\hat P(\Omega)|\psi\rangle
    }.
\]
Для нормированного состояния эта величина имеет смысл вероятности
получить при измерении значение наблюдаемой, принадлежащее множеству
$\Omega$.

Действительно,
\[
\begin{aligned}
    \|\hat P(\Omega)|\psi\rangle\|^2
    &=
    \langle\psi|
        \hat P^\dagger(\Omega)\hat P(\Omega)
    |\psi\rangle \\
    &=
    \langle\psi|\hat P(\Omega)|\psi\rangle,
\end{aligned}
\]
поскольку ортогональный проектор удовлетворяет
\[
    \hat P^\dagger(\Omega)=\hat P(\Omega),
    \qquad
    \hat P^2(\Omega)=\hat P(\Omega).
\]

Если
\[
    \Omega_1\cap\Omega_2=\varnothing,
\]
то
\[
    \hat P(\Omega_1\cup\Omega_2)
    =
    \hat P(\Omega_1)+\hat P(\Omega_2)
\]
и
\[
    \boxed{
        \hat P(\Omega_1)\hat P(\Omega_2)=0
    }.
\]
В общем случае проекторы, соответствующие одной спектральной мере,
коммутируют и
\[
    \boxed{
        \hat P(\Omega_1)\hat P(\Omega_2)
        =
        \hat P(\Omega_1\cap\Omega_2)
    }.
\]
Поэтому
\[
    \hat P(\Omega_1\cup\Omega_2)
    =
    \hat P(\Omega_1)+\hat P(\Omega_2)
    -
    \hat P(\Omega_1\cap\Omega_2).
\]

Для самосопряжённого оператора $\hat A$ существует соответствующая ему
спектральная мера
\[
    \{\hat P_A(\Omega)\}.
\]
Единичный оператор представляется как
\[
    \boxed{
        \hat I
        =
        \hat P_A(\mathbb R)
        =
        \int_{\mathbb R}d\hat P_A(\lambda)
    }.
\]


% =========================================================
\subsubsection{Спектральная теорема в общем случае}
% =========================================================

Спектральное разложение самосопряжённого оператора обобщает конечномерную
формулу
\[
    \hat A=\sum_n\lambda_n\hat P_n
\]
и имеет вид
\[
    \boxed{
        \hat A
        =
        \int_{\mathbb R}
            \lambda\,d\hat P_A(\lambda)
    }.
\]

Если оператор $\hat A$ неограничен, его область определения состоит из
тех состояний, для которых соответствующий интеграл сходится:
\[
    \boxed{
    \mathcal D(\hat A)
    =
    \left\{
        |\psi\rangle:
        \int_{\mathbb R}
            \lambda^2\,
            d\langle\psi|
                \hat P_A(\lambda)
            |\psi\rangle
        <\infty
    \right\}
    }.
\]
Таким образом, условие принадлежности области определения связано с
конечностью второго момента спектральной меры.

В дискретном случае эти формулы переходят в уже знакомые выражения
\[
    \hat I=\sum_n\hat P_n,
    \qquad
    \hat A=\sum_n\lambda_n\hat P_n.
\]

Для функции $f$ от самосопряжённого оператора спектральное разложение
задаёт
\[
    \boxed{
        f(\hat A)
        =
        \int_{\mathbb R}
            f(\lambda)\,
            d\hat P_A(\lambda)
    }.
\]
В частности, для резольвенты
\[
    \hat R_A(z)
    =
    (\hat A-z\hat I)^{-1}
\]
получаем
\[
    \boxed{
        \hat R_A(z)
        =
        \int_{\mathbb R}
        \frac{d\hat P_A(\lambda)}
             {\lambda-z}
    }.
\]

Точка $\lambda$ принадлежит спектру $\sigma(\hat A)$ тогда и только
тогда, когда любая её окрестность имеет ненулевой спектральный
проектор:
\[
    \boxed{
        \lambda\in\sigma(\hat A)
        \quad\Longleftrightarrow\quad
        \hat P_A((\lambda-\varepsilon,\lambda+\varepsilon))
        \neq0
        \quad
        \forall\,\varepsilon>0
    }.
\]

Из спектрального разложения следует также формула для среднего
значения:
\[
\begin{aligned}
    \langle\psi|\hat A|\psi\rangle
    &=
    \int_{\mathbb R}
        \lambda\,
        d\langle\psi|
            \hat P_A(\lambda)
        |\psi\rangle.
\end{aligned}
\]
Если обозначить
\[
    \boxed{
        M_\psi(\lambda)
        =
        \langle\psi|
            \hat P_A(\lambda)
        |\psi\rangle
    },
\]
то
\[
    \boxed{
        \langle\hat A\rangle_\psi
        =
        \int_{\mathbb R}
            \lambda\,dM_\psi(\lambda)
    }.
\]


% =========================================================
\subsubsection{Части спектральной меры}
% =========================================================

Спектральную меру можно представить как сумму трёх частей:
\[
    \boxed{
        M_\psi
        =
        M_\psi^{pp}
        +
        M_\psi^{ac}
        +
        M_\psi^{sc}
    }.
\]
Здесь $M_\psi^{pp}$ --- чисто точечная часть,
$M_\psi^{ac}$ --- абсолютно непрерывная часть, а
$M_\psi^{sc}$ --- сингулярно непрерывная часть.

Для абсолютно непрерывной части можно записать
\[
    dM_\psi^{ac}(\lambda)
    =
    \rho_\psi(\lambda)\,d\lambda,
\]
где $\rho_\psi(\lambda)$ играет роль плотности соответствующей
спектральной меры.

Этим трём типам меры соответствуют подпространства
\[
    \mathcal H^{pp},
    \qquad
    \mathcal H^{ac},
    \qquad
    \mathcal H^{sc},
\]
причём
\[
    \boxed{
        \mathcal H
        =
        \mathcal H^{pp}
        \oplus
        \mathcal H^{ac}
        \oplus
        \mathcal H^{sc}
    }.
\]
Каждое из этих подпространств инвариантно относительно $\hat A$:
\[
    \boxed{
        \hat A\mathcal H^i
        \subseteq
        \mathcal H^i,
        \qquad
        i=pp,ac,sc
    }.
\]


% =========================================================
\subsubsection{Нормальные и унитарные операторы}
% =========================================================

Спектральное разложение обобщается и на нормальные операторы.
Оператор $\hat N$ называется нормальным, если
\[
    \boxed{
        \hat N^\dagger\hat N
        =
        \hat N\hat N^\dagger
    }.
\]
В конечномерном случае для нормального оператора можно записать
\[
    \hat N
    =
    \sum_n\lambda_n\hat P_n,
    \qquad
    \lambda_n\in\mathbb C.
\]
В общем случае спектральная мера такого оператора задаётся уже на
подмножествах комплексной плоскости:
\[
    \hat P_N(\Omega),
    \qquad
    \Omega\subset\mathbb C.
\]

Нормальный оператор можно представить в форме
\[
    \boxed{
        \hat N=\hat A+i\hat B
    },
\]
где $\hat A$ и $\hat B$ --- самосопряжённые операторы.

Важным частным случаем нормального оператора является унитарный
оператор $\hat U$:
\[
    \boxed{
        \hat U^\dagger\hat U
        =
        \hat U\hat U^\dagger
        =
        \hat I
    }.
\]
Он сохраняет скалярное произведение:
\[
    \langle\hat U\psi|\hat U\psi\rangle
    =
    \langle\psi|
        \hat U^\dagger\hat U
    |\psi\rangle
    =
    \langle\psi|\psi\rangle.
\]
Следовательно, унитарное преобразование сохраняет норму состояния.

Если
\[
    \hat U|\psi\rangle
    =
    \lambda|\psi\rangle,
\]
то
\[
    |\lambda|=1,
\]
поэтому собственные значения унитарного оператора имеют вид
\[
    \boxed{
        \lambda=e^{i\alpha},
        \qquad
        \alpha\in\mathbb R
    }.
\]


% =========================================================
\subsubsection{Унитарная непрерывная группа}
% =========================================================

Рассмотрим однопараметрическое семейство унитарных операторов
$\hat U(t)$, $t\in\mathbb R$, для которого
\[
    \boxed{
        \hat U(0)=\hat I,
        \qquad
        \hat U(t)\hat U(s)
        =
        \hat U(t+s)
    }.
\]
Потребуем также непрерывности:
\[
    \boxed{
        \hat U(t)|\psi\rangle
        \longrightarrow
        \hat U(t_0)|\psi\rangle,
        \qquad
        t\to t_0
    }.
\]

Тогда существует самосопряжённый оператор $\hat A$, такой что
\[
    \boxed{
        \hat U(t)=e^{it\hat A}
    }.
\]
Оператор $\hat A$ называется генератором этой однопараметрической
унитарной группы.

Используя спектральное разложение $\hat A$, экспоненту можно записать
как
\[
    \boxed{
        \hat U(t)
        =
        \int_{\mathbb R}
            e^{it\lambda}\,
            d\hat P_A(\lambda)
    }.
\]


% =========================================================
\subsection{Совместно измеримые наблюдаемые. Соотношение неопределённостей}
% =========================================================

Рассмотрим две наблюдаемые, которым соответствуют самосопряжённые
операторы $\hat A$ и $\hat B$.

Пусть
\[
    \hat A|\psi\rangle
    =
    \lambda|\psi\rangle,
\]
то есть $|\psi\rangle$ принадлежит собственному подпространству
$\hat A$, соответствующему $\lambda$. Если
\[
    \hat A\hat B=\hat B\hat A,
\]
то
\[
    \hat A\hat B|\psi\rangle
    =
    \hat B\hat A|\psi\rangle
    =
    \lambda\hat B|\psi\rangle.
\]
Следовательно, $\hat B|\psi\rangle$ снова принадлежит тому же
собственному подпространству $\hat A$.

В общем случае совместная измеримость наблюдаемых выражается через
их спектральные проекторы:
\[
    \boxed{
        [\hat P_A(\Omega_1),\hat P_B(\Omega_2)]
        =
        0
        \qquad
        \forall\,\Omega_1,\Omega_2
    }.
\]
То есть спектральные проекторы двух совместно измеримых наблюдаемых
должны коммутировать.

В конечномерном случае это соответствует обычному условию
\[
    \boxed{
        \hat A\hat B=\hat B\hat A
    }.
\]
Для самосопряжённых операторов это условие удобно также выражать через
их унитарные группы:
\[
    \boxed{
        [e^{it\hat A},e^{is\hat B}]=0
        \qquad
        \forall\,t,s\in\mathbb R
    }.
\]

Набор наблюдаемых
\[
    \hat A_1,\ldots,\hat A_n
\]
называется полным набором совместно измеримых наблюдаемых, если они
совместно измеримы и любая наблюдаемая $\hat B$, совместно измеримая
со всеми $\hat A_i$, представляется как функция этого набора:
\[
    \boxed{
        \hat B
        =
        f(\hat A_1,\ldots,\hat A_n)
    }.
\]


% =========================================================
\subsubsection{Соотношение неопределённостей}
% =========================================================

Пусть $\hat A$ и $\hat B$ --- самосопряжённые операторы, а
$|\psi\rangle$ --- нормированное состояние, принадлежащее необходимым
областям определения. Обозначим средние значения
\[
    \overline A
    =
    \langle\psi|\hat A|\psi\rangle,
    \qquad
    \overline B
    =
    \langle\psi|\hat B|\psi\rangle.
\]
Введём операторы отклонений
\[
    \boxed{
        \hat A_0
        =
        \hat A-\overline A\,\hat I,
        \qquad
        \hat B_0
        =
        \hat B-\overline B\,\hat I
    }.
\]
Тогда дисперсии равны
\[
    \boxed{
        \operatorname{Var}A
        =
        \|\hat A_0|\psi\rangle\|^2,
        \qquad
        \operatorname{Var}B
        =
        \|\hat B_0|\psi\rangle\|^2
    }.
\]

Из неравенства Коши--Буняковского следует
\[
    \operatorname{Var}A\,
    \operatorname{Var}B
    \geq
    \left|
        \langle
            \hat A_0\psi
            |
            \hat B_0\psi
        \rangle
    \right|^2.
\]
Для комплексного числа можно отдельно выделить действительную и
мнимую части:
\begin{align*}
    \left|
        \langle\hat A_0\psi|\hat B_0\psi\rangle
    \right|^2
    &=
    \left|
        \frac{
            \langle\hat A_0\psi|\hat B_0\psi\rangle
            +
            \langle\hat B_0\psi|\hat A_0\psi\rangle
        }{2}
    \right|^2
    \\
    &\quad+
    \left|
        \frac{
            \langle\hat A_0\psi|\hat B_0\psi\rangle
            -
            \langle\hat B_0\psi|\hat A_0\psi\rangle
        }{2i}
    \right|^2.
\end{align*}
Отбрасывая первый, неотрицательный член, получаем
\[
    \operatorname{Var}A\,
    \operatorname{Var}B
    \geq
    \frac14
    \left|
        \langle\psi|
            [\hat A_0,\hat B_0]
        |\psi\rangle
    \right|^2.
\]
Но добавление к операторам постоянных, умноженных на единичный
оператор, не меняет коммутатор:
\[
    [\hat A_0,\hat B_0]
    =
    [\hat A,\hat B].
\]
Следовательно,
\[
    \boxed{
    \operatorname{Var}A\,
    \operatorname{Var}B
    \geq
    \frac14
    \left|
        \langle\psi|
            [\hat A,\hat B]
        |\psi\rangle
    \right|^2
    }.
\]

Если определить неопределённость наблюдаемой как
\[
    \sigma_A
    =
    \sqrt{\operatorname{Var}A},
    \qquad
    \sigma_B
    =
    \sqrt{\operatorname{Var}B},
\]
получаем соотношение неопределённостей
\[
    \boxed{
        \sigma_A\sigma_B
        \geq
        \frac12
        \left|
            \langle\psi|
                [\hat A,\hat B]
            |\psi\rangle
        \right|
    }.
\]


% =========================================================
\subsubsection{Коммутатор}
% =========================================================

Коммутатор двух операторов определяется как
\[
    \boxed{
        [\hat A,\hat B]
        =
        \hat A\hat B-\hat B\hat A
    }.
\]
Из определения непосредственно следует антисимметрия:
\[
    \boxed{
        [\hat A,\hat B]
        =
        -[\hat B,\hat A]
    }.
\]
Для произведения операторов выполняется правило
\[
    \boxed{
        [\hat A,\hat B\hat C]
        =
        [\hat A,\hat B]\hat C
        +
        \hat B[\hat A,\hat C]
    }.
\]
Кроме того, коммутатор удовлетворяет тождеству Якоби:
\[
    \boxed{
    [\hat A,[\hat B,\hat C]]
    +
    [\hat B,[\hat C,\hat A]]
    +
    [\hat C,[\hat A,\hat B]]
    =
    0
    }.
\]


% =========================================================
\subsubsection{Координата и импульс}
% =========================================================

В координатном представлении
\[
    \hat x=x,
    \qquad
    \hat p_x
    =
    -i\hbar\frac{d}{dx}.
\]
Рассмотрим их коммутатор, действующий на функцию $\psi(x)$:
\begin{align*}
    [\hat x,\hat p_x]\psi
    &=
    -i\hbar x\frac{d\psi}{dx}
    +
    i\hbar\frac{d}{dx}(x\psi) \\
    &=
    -i\hbar x\psi'
    +
    i\hbar\psi
    +
    i\hbar x\psi' \\
    &=
    i\hbar\psi.
\end{align*}
Следовательно,
\[
    \boxed{
        [\hat x,\hat p_x]
        =
        i\hbar\hat I
    }.
\]
Это каноническое коммутационное соотношение.

Поэтому для координаты и импульса соотношение неопределённостей принимает
вид
\[
    \boxed{
        \sigma_x\sigma_{p_x}
        \geq
        \frac{\hbar}{2}
    }.
\]

В трёхмерном случае
\[
    \boxed{
        [\hat x_\alpha,\hat p_\beta]
        =
        i\hbar\delta_{\alpha\beta}\hat I
    }.
\]

% =========================================================
\section{Лекции 15--16}
% =========================================================

\subsection{Коммутаторы и соотношение неопределённостей}

Для двух наблюдаемых, которым соответствуют самосопряжённые операторы
$\hat A$ и $\hat B$, соотношение неопределённостей имеет вид
\[
    \boxed{
        \Delta A\,\Delta B
        \geq
        \frac12
        \left|
            \langle\psi|[\hat A,\hat B]|\psi\rangle
        \right|
    }.
\]
Здесь
\[
    [\hat A,\hat B]=\hat A\hat B-\hat B\hat A
\]
--- коммутатор операторов.

Для координаты и импульса выполняются канонические коммутационные
соотношения
\[
    \boxed{
        [\hat x_\alpha,\hat p_\beta]
        =
        i\hbar\delta_{\alpha\beta}\hat I
    }.
\]

Если $F(\hat p)$ является достаточно регулярной функцией импульса, то
\[
    \boxed{
        [\hat x,F(\hat p)]
        =
        i\hbar\frac{\partial F}{\partial p}
    }.
\]
Аналогично, для функции координаты $G(\hat x)$
\[
    \boxed{
        [\hat p,G(\hat x)]
        =
        -i\hbar\frac{\partial G}{\partial x}
    }.
\]
Например, при
\[
    F(p)=p^2
\]
получаем
\[
    \boxed{
        [\hat x,\hat p^2]
        =
        2i\hbar\hat p
    }.
\]

\subsubsection{Коммутационные соотношения для момента импульса}

Введём безразмерный оператор орбитального момента импульса
\[
    \boxed{
        \hat{\vec L}
        =
        \frac{1}{\hbar}
        [\hat{\vec r}\times\hat{\vec p}]
    }.
\]
Его компоненты удовлетворяют следующим коммутационным соотношениям:
\[
    \boxed{
        [\hat x_\alpha,\hat L_\beta]
        =
        i\varepsilon_{\alpha\beta\gamma}\hat x_\gamma
    },
\]
\[
    \boxed{
        [\hat p_\alpha,\hat L_\beta]
        =
        i\varepsilon_{\alpha\beta\gamma}\hat p_\gamma
    },
\]
\[
    \boxed{
        [\hat L_\alpha,\hat L_\beta]
        =
        i\varepsilon_{\alpha\beta\gamma}\hat L_\gamma
    }.
\]
В частности,
\[
    [\hat L_x,\hat L_y]
    =
    i\hat L_z.
\]

Если
\[
    \hat L_z|m\rangle=m|m\rangle,
\]
то
\[
    \langle m|\hat L_z|m\rangle=m.
\]

Рассмотрим угловую координату $\varphi\in[0,2\pi)$:
\[
    \hat L_z
    =
    -i\frac{\partial}{\partial\varphi}.
\]
Её собственные функции имеют вид
\[
    \boxed{
        \psi_m(\varphi)
        =
        \frac{1}{\sqrt{2\pi}}e^{im\varphi},
        \qquad
        m\in\mathbb Z
    }.
\]
Формально
\[
    [\varphi,\hat L_z]=i\hat I.
\]
Однако попытка вычислить среднее значение этого коммутатора в
собственном состоянии $\hat L_z$ приводит к
\begin{align*}
    \langle m|[\varphi,\hat L_z]|m\rangle
    &=
    \langle m|\varphi\hat L_z|m\rangle
    -
    \langle m|\hat L_z\varphi|m\rangle
    \\
    &=
    m\langle m|\varphi|m\rangle
    -
    m\langle m|\varphi|m\rangle
    =0,
\end{align*}
что на первый взгляд противоречит формальному соотношению
$[\varphi,\hat L_z]=i\hat I$.

Причина заключается в областях определения операторов.

Для оператора $\hat L_z$ естественны периодические граничные условия
\[
    \psi(0)=\psi(2\pi),
\]
то есть
\[
    \mathcal D(\hat L_z)
    =
    \left\{
        \psi\in AC[0,2\pi]:
        \psi(0)=\psi(2\pi)
    \right\}.
\]
Но для произведения $\hat L_z\varphi$ требуется, чтобы функция
$\varphi\psi(\varphi)$ также принадлежала области определения
$\hat L_z$:
\[
    \left.
        \varphi\psi(\varphi)
    \right|_{\varphi=0}
    =
    \left.
        \varphi\psi(\varphi)
    \right|_{\varphi=2\pi}.
\]
Отсюда
\[
    0
    =
    2\pi\psi(2\pi),
\]
то есть дополнительно требуется
\[
    \psi(2\pi)=0.
\]
Собственные функции $e^{im\varphi}$ этому условию не удовлетворяют.
Следовательно, в данном случае нельзя непосредственно применять
формулу с коммутатором к состоянию $|m\rangle$.

Более общая форма соотношения неопределённостей, не требующая
существования обоих произведений $\hat A\hat B$ и $\hat B\hat A$,
имеет вид
\[
    \boxed{
    \Delta A\,\Delta B
    \geq
    \frac12
    \left|
        \langle\hat A\psi|\hat B\psi\rangle
        -
        \langle\hat B\psi|\hat A\psi\rangle
    \right|
    },
\]
где достаточно
\[
    |\psi\rangle
    \in
    \mathcal D(\hat A)\cap\mathcal D(\hat B).
\]

Применим эту формулу к
\[
    \hat A=\varphi,
    \qquad
    \hat B=-i\frac{\partial}{\partial\varphi}.
\]
Тогда
\begin{align*}
    &
    \langle\varphi\psi|\hat L_z\psi\rangle
    -
    \langle\hat L_z\psi|\varphi\psi\rangle
    \\
    &=
    -i
    \int_0^{2\pi}
        \varphi\psi^*(\varphi)
        \frac{\partial\psi}{\partial\varphi}\,d\varphi
    -
    i
    \int_0^{2\pi}
        \varphi
        \frac{\partial\psi^*}{\partial\varphi}
        \psi(\varphi)\,d\varphi
    \\
    &=
    -i
    \int_0^{2\pi}
        \varphi
        \frac{\partial}{\partial\varphi}
        |\psi(\varphi)|^2\,d\varphi
    \\
    &=
    -i
    \left.
        \varphi|\psi(\varphi)|^2
    \right|_0^{2\pi}
    +
    i
    \int_0^{2\pi}
        |\psi(\varphi)|^2\,d\varphi.
\end{align*}
Для нормированной функции
\[
    \int_0^{2\pi}|\psi(\varphi)|^2\,d\varphi=1,
\]
поэтому
\[
    \boxed{
    \langle\varphi\psi|\hat L_z\psi\rangle
    -
    \langle\hat L_z\psi|\varphi\psi\rangle
    =
    i\left(
        1-2\pi|\psi(2\pi)|^2
    \right)
    }.
\]
Для собственного состояния
\[
    \psi_m(\varphi)
    =
    \frac{1}{\sqrt{2\pi}}e^{im\varphi}
\]
правая часть действительно обращается в нуль.


% =========================================================
\subsection{Представления в квантовой механике}
% =========================================================

Выбор представления в квантовой механике соответствует выбору базиса
в гильбертовом пространстве состояний.

Рассмотрим наблюдаемую $\hat A$ с простым спектром. Состоянию
$|\psi\rangle$ можно сопоставить функцию
\[
    \boxed{
        \psi(\lambda)
        =
        \langle\lambda|\psi\rangle
    }.
\]
В таком представлении
\[
    \mathcal H
    \simeq
    L_2(\mathbb R,d\mu(\lambda)),
\]
а норма состояния принимает вид
\[
    \boxed{
        \langle\psi|\psi\rangle
        =
        \int_{\mathbb R}
            |\psi(\lambda)|^2\,d\mu(\lambda)
    }.
\]
Формально это можно записать как
\[
    \langle\psi|\psi\rangle
    =
    \int_{\mathbb R}
        \langle\psi|\lambda\rangle
        \langle\lambda|\psi\rangle
        \,d\mu(\lambda).
\]

Пусть
\[
    |\eta\rangle=\hat A|\psi\rangle.
\]
Тогда в спектральном представлении
\[
    \eta(\lambda)
    =
    \langle\lambda|\hat A|\psi\rangle
    =
    \lambda\psi(\lambda).
\]
Таким образом, оператор $\hat A$ в своём собственном представлении
действует как умножение на спектральную переменную:
\[
    \boxed{
        \hat A
        \longrightarrow
        \lambda
    }.
\]

Аналогично,
\begin{align*}
    \langle\psi|f(\hat A)|\psi\rangle
    &=
    \int_{\mathbb R}
        \psi^*(\lambda)
        f(\lambda)
        \psi(\lambda)
        \,d\mu(\lambda)
    \\
    &=
    \int_{\mathbb R}
        \langle\psi|\lambda\rangle
        f(\lambda)
        \langle\lambda|\psi\rangle
        \,d\mu(\lambda).
\end{align*}

Для неограниченного оператора область определения задаётся условием
\[
    \boxed{
    \mathcal D(\hat A)
    =
    \left\{
        \psi:
        \int_{\mathbb R}
            \lambda^2
            |\psi(\lambda)|^2
            \,d\mu(\lambda)
        <\infty
    \right\}
    }.
\]

\subsubsection{Дискретный спектр}

Если спектр дискретен и
\[
    \hat A|\lambda_n\rangle
    =
    \lambda_n|\lambda_n\rangle,
\]
то при простом спектре собственные подпространства одномерны:
\[
    \dim\operatorname{Ran}\hat P_n=1.
\]
Разложение состояния имеет вид
\[
    \boxed{
        |\psi\rangle
        =
        \sum_{n=1}^{N}
        |\lambda_n\rangle
        \langle\lambda_n|\psi\rangle
    }.
\]
В собственном базисе оператор представляется диагональной матрицей
\[
    \hat A
    \longleftrightarrow
    \begin{pmatrix}
        \lambda_1 & & 0\\
        & \lambda_2 & \\
        0 & & \ddots
    \end{pmatrix}.
\]
Среднее значение равно
\[
    \boxed{
        \langle\psi|\hat A|\psi\rangle
        =
        \sum_n
            \lambda_n
            |\langle\lambda_n|\psi\rangle|^2
    }.
\]

\subsubsection{Набор совместимых наблюдаемых}

Если одного квантового числа недостаточно для однозначной нумерации
состояний, используется набор наблюдаемых
\[
    \vec\lambda
    =
    (\lambda_1,\ldots,\lambda_k).
\]
Тогда
\[
    \mathcal H
    \simeq
    L_2(\mathbb R^k,d\mu(\vec\lambda)),
\]
а
\[
    |\psi\rangle
    \longrightarrow
    \psi(\vec\lambda)
    =
    \langle\vec\lambda|\psi\rangle.
\]
Операторы соответствующих наблюдаемых действуют как
\[
    \boxed{
        \hat A_j\psi(\vec\lambda)
        =
        \lambda_j\psi(\vec\lambda)
    }.
\]
Мера $d\mu(\vec\lambda)$ в общем случае не обязана раскладываться в
произведение одномерных мер:
\[
    d\mu(\vec\lambda)
    \neq
    d\mu_1(\lambda_1)
    \cdots
    d\mu_k(\lambda_k).
\]

Примером нескольких квантовых чисел являются состояния атома
водорода. Для связанных состояний при фиксированном главном квантовом
числе $n$ используются также
\[
    l=0,\ldots,n-1,
    \qquad
    m=-l,\ldots,l.
\]


% =========================================================
\subsubsection{Дискретная и непрерывная части спектра}
% =========================================================

Спектральную меру можно представить в виде
\[
    \boxed{
        d\mu(\lambda)
        =
        d\mu_{pp}(\lambda)
        +
        d\mu_{ac}(\lambda)
        +
        d\mu_{sc}(\lambda)
    }.
\]
Чисто точечная часть имеет вид
\[
    d\mu_{pp}(\lambda)
    =
    \sum_n
        p_n\,
        \delta(\lambda-\lambda_n)\,d\lambda,
\]
а абсолютно непрерывную часть можно записать как
\[
    d\mu_{ac}(\lambda)
    =
    \nu(\lambda)\,d\lambda.
\]

Если сингулярно непрерывная часть отсутствует,
\[
    d\mu_{sc}=0,
\]
то условие полноты можно записать в виде
\[
\boxed{
\langle\psi|\psi\rangle
=
\sum_n
    \langle\psi|\lambda_n\rangle
    \langle\lambda_n|\psi\rangle
+
\int_{\sigma_{\mathrm{cont}}}
    d\lambda\,
    \langle\psi|\lambda\rangle
    \langle\lambda|\psi\rangle
}.
\]

Формально единичный оператор имеет вид
\[
    \boxed{
        \hat I
        =
        \int
            |\lambda\rangle
            d\mu(\lambda)
            \langle\lambda|
    }.
\]
Для обобщённых собственных векторов
\[
    \langle\lambda|\hat A|\psi\rangle
    =
    \lambda\langle\lambda|\psi\rangle.
\]
Из самосопряжённости также формально следует
\[
    \langle\psi|\hat A|\lambda\rangle
    =
    \lambda\langle\psi|\lambda\rangle,
\]
то есть
\[
    \boxed{
        \hat A|\lambda\rangle
        =
        \lambda|\lambda\rangle
    }.
\]
При непрерывном спектре такие $|\lambda\rangle$ обычно не принадлежат
самому гильбертову пространству и рассматриваются как обобщённые
собственные векторы.


% =========================================================
\subsubsection{Оснащённое гильбертово пространство}
% =========================================================

Для аккуратного описания обобщённых собственных векторов используется
цепочка пространств
\[
    \boxed{
        \mathcal H_+
        \subset
        \mathcal H
        \simeq
        \mathcal H^*
        \subset
        \mathcal H_-
    }.
\]
Пространство $\mathcal H_+$ выбирается достаточно ``хорошим'', например
состоящим из достаточно гладких функций. Его элементы можно считать
физически допустимыми пробными состояниями.

Обобщённый собственный вектор действует на
$|\psi\rangle\in\mathcal H_+$ как линейный функционал
\[
    \lambda(\psi):
    \mathcal H_+\to\mathbb C.
\]
В обозначениях Дирака этот функционал записывается как
\[
    \langle\lambda|\psi\rangle.
\]
Сами обобщённые собственные векторы относятся к расширенному
пространству $\mathcal H_-$.


% =========================================================
\subsubsection{Импульсное представление}
% =========================================================

Для свободной частицы
\[
    \hat p|p\rangle
    =
    p|p\rangle.
\]
В координатном представлении
\[
    \boxed{
        \langle x|p\rangle
        =
        \frac{1}{\sqrt{2\pi\hbar}}
        e^{ipx/\hbar}
    }.
\]
Плоская волна не принадлежит $L_2(\mathbb R)$, поэтому $|p\rangle$
является обобщённым собственным состоянием.

Импульсная волновая функция определяется как
\[
    \boxed{
    \langle p|\psi\rangle
    =
    \psi(p)
    =
    \frac{1}{\sqrt{2\pi\hbar}}
    \int_{-\infty}^{+\infty}
        e^{-ipx/\hbar}\psi(x)\,dx
    }.
\]
Таким образом, переход
\[
    \psi(x)\longrightarrow\psi(p)
\]
является преобразованием Фурье.


% =========================================================
\subsubsection{Матричные элементы и интегральное ядро оператора}
% =========================================================

Пусть $\hat B$ --- некоторый оператор. В спектральном представлении
оператора $\hat A$
\begin{align*}
    \langle\psi|\hat B|\psi\rangle
    &=
    \iint
        d\mu(\lambda')
        d\mu(\lambda)\,
        \langle\psi|\lambda'\rangle
        \langle\lambda'|\hat B|\lambda\rangle
        \langle\lambda|\psi\rangle
    \\
    &=
    \iint
        d\mu(\lambda')
        d\mu(\lambda)\,
        \psi^*(\lambda')
        B(\lambda',\lambda)
        \psi(\lambda),
\end{align*}
где
\[
    \boxed{
        B(\lambda',\lambda)
        =
        \langle\lambda'|\hat B|\lambda\rangle
    }
\]
--- интегральное ядро оператора.

Функция самого оператора $\hat A$ имеет диагональное ядро:
\[
    \boxed{
        f(\hat A)
        \longrightarrow
        f(\lambda)\,
        \delta_\mu(\lambda-\lambda')
    }.
\]
Для единичного оператора
\[
    \boxed{
        \hat I
        \longrightarrow
        \delta_\mu(\lambda-\lambda')
    }.
\]
Здесь $\delta_\mu$ понимается как $\delta$-функция относительно
спектральной меры, то есть
\[
    \int
        \delta_\mu(\lambda-\lambda')
        f(\lambda')\,d\mu(\lambda')
    =
    f(\lambda).
\]


% =========================================================
\subsubsection{Координатное представление}
% =========================================================

Для оператора координаты
\[
    \hat x|\psi\rangle
    \longrightarrow
    \langle x|\hat x|\psi\rangle
    =
    x\psi(x).
\]
Обобщённые собственные векторы координаты удовлетворяют
\[
    \boxed{
        \hat x|x_0\rangle
        =
        x_0|x_0\rangle
    }.
\]
Используя полноту координатного базиса,
\[
    \langle\psi|x_0\rangle
    =
    \int_{-\infty}^{+\infty}
        dx\,
        \langle\psi|x\rangle
        \langle x|x_0\rangle.
\]
Левая часть равна $\psi^*(x_0)$, поэтому для произвольной $\psi$
необходимо
\[
    \boxed{
        \langle x|x_0\rangle
        =
        \delta(x-x_0)
    }.
\]
То есть $|x_0\rangle$ является обобщённым вектором, соответствующим
$\delta$-функции.

Отсюда
\[
    \boxed{
        \langle x|\hat x|x'\rangle
        =
        x\,\delta(x-x')
    }.
\]

Для оператора импульса
\[
    \hat p
    =
    -i\hbar\frac{\partial}{\partial x}
\]
имеем
\[
    \langle x|\hat p|\psi\rangle
    =
    -i\hbar
    \frac{\partial}{\partial x}
    \psi(x).
\]
Вставляя разложение единицы,
\[
    \langle x|\hat p|\psi\rangle
    =
    \int dx'\,
        \langle x|\hat p|x'\rangle
        \langle x'|\psi\rangle,
\]
и сравнивая с
\[
    -i\hbar
    \frac{\partial}{\partial x}
    \int dx'\,
        \delta(x-x')
        \langle x'|\psi\rangle,
\]
получаем
\[
    \boxed{
        \langle x|\hat p|x'\rangle
        =
        -i\hbar
        \frac{\partial}{\partial x}
        \delta(x-x')
        =
        i\hbar
        \frac{\partial}{\partial x'}
        \delta(x-x')
    }.
\]


% =========================================================
\subsubsection{Смена представления}
% =========================================================

Пусть имеются два полных набора обобщённых базисных состояний
\[
    \{|\xi\rangle\},
    \qquad
    \{|\eta\rangle\}.
\]
Для них
\[
    \hat I
    =
    \int d\xi\,|\xi\rangle\langle\xi|
    =
    \int d\eta\,|\eta\rangle\langle\eta|.
\]
Волновые функции в двух представлениях связаны формулами
\[
    \boxed{
    \psi(\xi)
    =
    \langle\xi|\psi\rangle
    =
    \int d\eta\,
        \langle\xi|\eta\rangle
        \psi(\eta)
    },
\]
\[
    \boxed{
    \psi(\eta)
    =
    \int d\xi\,
        \langle\eta|\xi\rangle
        \psi(\xi)
    }.
\]

Из полноты базисов следуют условия
\[
    \boxed{
    \int d\eta\,
        \langle\xi|\eta\rangle
        \langle\eta|\xi'\rangle
    =
    \delta(\xi-\xi')
    },
\]
\[
    \boxed{
    \int d\xi\,
        \langle\eta|\xi\rangle
        \langle\xi|\eta'\rangle
    =
    \delta(\eta-\eta')
    }.
\]

Преобразование между двумя представлениями является унитарным:
\[
    U:L_2^{(\xi)}\longrightarrow L_2^{(\eta)},
\]
\[
    \boxed{
        U^\dagger U=\hat I_\xi,
        \qquad
        UU^\dagger=\hat I_\eta
    }.
\]

Для координатного и импульсного представлений ядро этого преобразования
равно
\[
    \boxed{
        \langle x|p\rangle
        =
        \frac{1}{\sqrt{2\pi\hbar}}
        e^{ipx/\hbar}
    }.
\]


% =========================================================
\subsection{Линейный гармонический осциллятор в матричной формулировке}
% =========================================================

Гамильтониан линейного гармонического осциллятора равен
\[
    \boxed{
        \hat H
        =
        \frac{\hat p^2}{2m}
        +
        \frac{m\omega^2}{2}\hat x^2
    }.
\]
Введём безразмерные операторы $\hat\xi$ и $\hat\eta$:
\[
    \boxed{
        \hat x
        =
        \sqrt{\frac{\hbar}{m\omega}}\,
        \hat\xi,
        \qquad
        \hat p
        =
        \sqrt{m\hbar\omega}\,
        \hat\eta
    }.
\]
Из
\[
    [\hat x,\hat p]=i\hbar\hat I
\]
следует
\[
    \boxed{
        [\hat\xi,\hat\eta]
        =
        i\hat I
    }.
\]
Гамильтониан принимает вид
\[
    \boxed{
        \hat H
        =
        \frac{\hbar\omega}{2}
        \left(
            \hat\eta^2+\hat\xi^2
        \right)
    }.
\]

\subsubsection{Операторы рождения и уничтожения}

Определим
\[
    \boxed{
        \hat a
        =
        \frac{1}{\sqrt2}
        (\hat\xi+i\hat\eta)
    },
\]
\[
    \boxed{
        \hat a^\dagger
        =
        \frac{1}{\sqrt2}
        (\hat\xi-i\hat\eta)
    }.
\]
Обратные формулы:
\[
    \boxed{
        \hat\xi
        =
        \frac{\hat a+\hat a^\dagger}{\sqrt2},
        \qquad
        \hat\eta
        =
        \frac{\hat a-\hat a^\dagger}{i\sqrt2}
    }.
\]

Из канонического коммутационного соотношения получаем
\[
    \boxed{
        [\hat a,\hat a^\dagger]=\hat I
    },
\]
то есть
\[
    \hat a\hat a^\dagger
    =
    \hat a^\dagger\hat a+\hat I.
\]

Поэтому
\begin{align*}
    \hat H
    &=
    \frac{\hbar\omega}{2}
    \left(
        \hat a\hat a^\dagger
        +
        \hat a^\dagger\hat a
    \right)
    \\
    &=
    \boxed{
        \hbar\omega
        \left(
            \hat n+\frac12
        \right)
    },
\end{align*}
где
\[
    \boxed{
        \hat n=\hat a^\dagger\hat a
    }
\]
--- оператор числа возбуждений.

Оператор $\hat n$ является неотрицательным:
\[
    \langle\psi|\hat n|\psi\rangle
    =
    \|\hat a|\psi\rangle\|^2
    \geq0.
\]
Пусть
\[
    \hat n|\lambda\rangle
    =
    \lambda|\lambda\rangle,
    \qquad
    \lambda\geq0.
\]

Рассмотрим действие $\hat a$:
\[
\begin{aligned}
    \lambda\hat a|\lambda\rangle
    &=
    \hat a\hat n|\lambda\rangle
    =
    \hat a\hat a^\dagger\hat a|\lambda\rangle
    \\
    &=
    (\hat n+\hat I)
    \hat a|\lambda\rangle.
\end{aligned}
\]
Следовательно,
\[
    \boxed{
        \hat n\hat a|\lambda\rangle
        =
        (\lambda-1)\hat a|\lambda\rangle
    }.
\]
То есть $\hat a$ уменьшает собственное значение $\hat n$ на единицу.

Аналогично,
\[
    \boxed{
        \hat n\hat a^\dagger|\lambda\rangle
        =
        (\lambda+1)
        \hat a^\dagger|\lambda\rangle
    }.
\]
Поэтому $\hat a^\dagger$ повышает собственное значение на единицу.

Оператор $\hat a$ называется оператором уничтожения, а
$\hat a^\dagger$ --- оператором рождения.


% =========================================================
\subsubsection{Спектр оператора числа возбуждений}
% =========================================================

Поскольку $\hat n$ неотрицателен, последовательное применение
$\hat a$ не может бесконечно уменьшать собственное значение.
Следовательно, существует состояние $|0\rangle$, для которого
\[
    \boxed{
        \hat a|0\rangle=0
    }.
\]
Тогда
\[
    \hat n|0\rangle
    =
    \hat a^\dagger\hat a|0\rangle
    =
    0.
\]
Это основное состояние осциллятора.

Таким образом, собственные значения $\hat n$ равны
\[
    \boxed{
        n=0,1,2,\ldots
    }.
\]

Пусть
\[
    \hat a|n\rangle
    =
    c_n|n-1\rangle.
\]
Из
\[
    n
    =
    \langle n|
        \hat a^\dagger\hat a
    |n\rangle
    =
    |c_n|^2
\]
получаем
\[
    \boxed{
        \hat a|n\rangle
        =
        \sqrt n\,|n-1\rangle
    }.
\]
Аналогично,
\[
    \boxed{
        \hat a^\dagger|n\rangle
        =
        \sqrt{n+1}\,|n+1\rangle
    }.
\]

Поэтому
\[
    \boxed{
        |n+1\rangle
        =
        \frac{\hat a^\dagger}{\sqrt{n+1}}
        |n\rangle
    },
\]
и все состояния строятся из основного:
\[
    \boxed{
        |n\rangle
        =
        \frac{(\hat a^\dagger)^n}{\sqrt{n!}}
        |0\rangle
    }.
\]

Матричные представления операторов $\hat a$ и $\hat a^\dagger$ в
базисе $\{|n\rangle\}$ имеют вид
\[
    \hat a
    =
    \begin{pmatrix}
        0 & \sqrt1 & 0 & 0 & \cdots\\
        0 & 0 & \sqrt2 & 0 & \cdots\\
        0 & 0 & 0 & \sqrt3 & \cdots\\
        \vdots & \vdots & \vdots & \ddots & \ddots
    \end{pmatrix},
\]
\[
    \hat a^\dagger
    =
    \begin{pmatrix}
        0 & 0 & 0 & \cdots\\
        \sqrt1 & 0 & 0 & \cdots\\
        0 & \sqrt2 & 0 & \cdots\\
        0 & 0 & \sqrt3 & \ddots\\
        \vdots & \vdots & \ddots & \ddots
    \end{pmatrix}.
\]

Поскольку $\hat x$ и $\hat p$ выражаются через $\hat a$ и
$\hat a^\dagger$, их диагональные матричные элементы равны нулю:
\[
    \boxed{
        \langle n|\hat x|n\rangle=0,
        \qquad
        \langle n|\hat p|n\rangle=0
    }.
\]

Энергетический спектр имеет вид
\[
    \boxed{
        E_n
        =
        \hbar\omega
        \left(
            n+\frac12
        \right)
    }.
\]


% =========================================================
\subsubsection{Координатное представление}
% =========================================================

В безразмерной координате $\xi$
\[
    \hat\xi=\xi,
    \qquad
    \hat\eta
    =
    -i\frac{\partial}{\partial\xi}.
\]
Поэтому
\[
    \hat a
    =
    \frac{1}{\sqrt2}
    \left(
        \xi+\frac{\partial}{\partial\xi}
    \right).
\]

Пусть
\[
    \psi_0(\xi)
    =
    \langle\xi|0\rangle.
\]
Условие
\[
    \hat a|0\rangle=0
\]
даёт
\[
    \left(
        \xi+\frac{\partial}{\partial\xi}
    \right)
    \psi_0(\xi)
    =
    0.
\]
Отсюда
\[
    \psi_0(\xi)
    =
    C e^{-\xi^2/2}.
\]

Так как
\[
    x
    =
    \sqrt{\frac{\hbar}{m\omega}}\,\xi,
    \qquad
    dx
    =
    \sqrt{\frac{\hbar}{m\omega}}\,d\xi,
\]
условие нормировки
\[
    \int_{-\infty}^{+\infty}
        |\psi_0(x)|^2\,dx
    =
    1
\]
даёт
\[
    \boxed{
        C
        =
        \left(
            \frac{m\omega}{\pi\hbar}
        \right)^{1/4}
    }.
\]

Оператор рождения в координатном представлении равен
\[
    \boxed{
        \hat a^\dagger
        =
        \frac{1}{\sqrt2}
        \left(
            \xi-\frac{\partial}{\partial\xi}
        \right)
    }.
\]
Следовательно,
\[
\begin{aligned}
    \langle\xi|n\rangle
    &=
    \left(
        \frac{m\omega}{\pi\hbar}
    \right)^{1/4}
    \frac{1}{\sqrt{2^n n!}}
    \left(
        \xi-\frac{\partial}{\partial\xi}
    \right)^n
    e^{-\xi^2/2}.
\end{aligned}
\]

Используя тождество
\[
    \left(
        \xi-\frac{d}{d\xi}
    \right)
    \left[
        e^{-\xi^2/2}F(\xi)
    \right]
    =
    -e^{\xi^2/2}
    \frac{d}{d\xi}
    \left[
        e^{-\xi^2}F(\xi)
    \right],
\]
получаем
\[
    \boxed{
    \langle\xi|n\rangle
    =
    \left(
        \frac{m\omega}{\pi\hbar}
    \right)^{1/4}
    \frac{e^{-\xi^2/2}}{\sqrt{2^n n!}}
    H_n(\xi)
    },
\]
где
\[
    \boxed{
        H_n(\xi)
        =
        (-1)^n
        e^{\xi^2}
        \frac{d^n}{d\xi^n}
        e^{-\xi^2}
    }
\]
--- полином Эрмита.


% =========================================================
\subsubsection{Когерентные состояния}
% =========================================================

Когерентным состоянием называется собственное состояние оператора
уничтожения:
\[
    \boxed{
        \hat a|\alpha\rangle
        =
        \alpha|\alpha\rangle
    },
    \qquad
    \alpha\in\mathbb C.
\]

Разложим его по состояниям числа возбуждений:
\[
    |\alpha\rangle
    =
    \sum_{n=0}^{\infty}
        |n\rangle\langle n|\alpha\rangle.
\]
Используя
\[
    |n\rangle
    =
    \frac{(\hat a^\dagger)^n}{\sqrt{n!}}|0\rangle,
\]
получаем
\[
    |\alpha\rangle
    =
    \sum_{n=0}^{\infty}
    \frac{(\hat a^\dagger)^n}{\sqrt{n!}}
    |0\rangle
    \langle n|\alpha\rangle.
\]

Из уравнения
\[
    \hat a|\alpha\rangle
    =
    \alpha|\alpha\rangle
\]
следует рекуррентное соотношение
\[
    \boxed{
        \langle n+1|\alpha\rangle
        =
        \frac{\alpha}{\sqrt{n+1}}
        \langle n|\alpha\rangle
    }.
\]
Поэтому
\[
    \boxed{
        \langle n|\alpha\rangle
        =
        \frac{\alpha^n}{\sqrt{n!}}
        \langle0|\alpha\rangle
    }.
\]

Следовательно,
\begin{align*}
    |\alpha\rangle
    &=
    C
    \sum_{n=0}^{\infty}
        \frac{\alpha^n}{\sqrt{n!}}
        |n\rangle
    \\
    &=
    C
    \sum_{n=0}^{\infty}
        \frac{\alpha^n}{n!}
        (\hat a^\dagger)^n|0\rangle
    \\
    &=
    C
    e^{\alpha\hat a^\dagger}
    |0\rangle.
\end{align*}

Из нормировки
\[
    1
    =
    \langle\alpha|\alpha\rangle
    =
    |C|^2
    \sum_{n=0}^{\infty}
        \frac{|\alpha|^{2n}}{n!}
    =
    |C|^2e^{|\alpha|^2}
\]
получаем
\[
    C=e^{-|\alpha|^2/2}
\]
с точностью до общего фазового множителя.

Итак,
\[
    \boxed{
        |\alpha\rangle
        =
        e^{-|\alpha|^2/2}
        e^{\alpha\hat a^\dagger}
        |0\rangle
    }.
\]

Когерентное состояние является суперпозицией состояний
$|n\rangle$. Вероятность обнаружить $n$ квантов равна
\[
    \boxed{
        |\langle n|\alpha\rangle|^2
        =
        \frac{|\alpha|^{2n}}{n!}
        e^{-|\alpha|^2}
    }.
\]
Это распределение Пуассона с параметром $|\alpha|^2$.

Для когерентных состояний соотношение неопределённостей координаты и
импульса достигает минимально возможного значения, то есть в
соотношении неопределённостей реализуется равенство.

% =========================================================
\subsubsection{Упражнение: соотношение неопределённостей для когерентного состояния}
% =========================================================

Покажем, что когерентное состояние
\[
    \hat a|\alpha\rangle=\alpha|\alpha\rangle
\]
реализует равенство в соотношении неопределённостей для координаты и
импульса:
\[
    \Delta x\,\Delta p=\frac{\hbar}{2}.
\]

Из определения операторов рождения и уничтожения
\[
    \hat a
    =
    \frac{1}{\sqrt2}
    \left(
        \hat\xi+i\hat\eta
    \right),
    \qquad
    \hat a^\dagger
    =
    \frac{1}{\sqrt2}
    \left(
        \hat\xi-i\hat\eta
    \right)
\]
следует
\[
    \hat\xi
    =
    \frac{\hat a+\hat a^\dagger}{\sqrt2},
    \qquad
    \hat\eta
    =
    \frac{\hat a-\hat a^\dagger}{i\sqrt2}.
\]
Поэтому
\[
    \boxed{
    \hat x
    =
    \sqrt{\frac{\hbar}{2m\omega}}
    \left(
        \hat a+\hat a^\dagger
    \right)
    },
\]
\[
    \boxed{
    \hat p
    =
    -i
    \sqrt{\frac{m\hbar\omega}{2}}
    \left(
        \hat a-\hat a^\dagger
    \right)
    }.
\]

Так как
\[
    \hat a|\alpha\rangle
    =
    \alpha|\alpha\rangle,
\]
то
\[
    \langle\alpha|\hat a|\alpha\rangle
    =
    \alpha,
    \qquad
    \langle\alpha|\hat a^\dagger|\alpha\rangle
    =
    \alpha^*.
\]
Кроме того,
\[
    \langle\hat a^\dagger\hat a\rangle
    =
    |\alpha|^2,
\]
а из коммутационного соотношения
\[
    [\hat a,\hat a^\dagger]=\hat I
\]
следует
\[
    \langle\hat a\hat a^\dagger\rangle
    =
    |\alpha|^2+1.
\]
Также
\[
    \langle\hat a^2\rangle=\alpha^2,
    \qquad
    \langle(\hat a^\dagger)^2\rangle=(\alpha^*)^2.
\]

Среднее значение координаты равно
\[
    \boxed{
    \langle\hat x\rangle
    =
    \sqrt{\frac{\hbar}{2m\omega}}
    \left(
        \alpha+\alpha^*
    \right)
    }.
\]
Для квадрата координаты
\begin{align*}
    \langle\hat x^2\rangle
    &=
    \frac{\hbar}{2m\omega}
    \left\langle
        \hat a^2
        +
        \hat a\hat a^\dagger
        +
        \hat a^\dagger\hat a
        +
        (\hat a^\dagger)^2
    \right\rangle
    \\
    &=
    \frac{\hbar}{2m\omega}
    \left[
        \alpha^2
        +
        (\alpha^*)^2
        +
        2|\alpha|^2
        +
        1
    \right].
\end{align*}
С другой стороны,
\[
    \langle\hat x\rangle^2
    =
    \frac{\hbar}{2m\omega}
    \left[
        \alpha^2
        +
        (\alpha^*)^2
        +
        2|\alpha|^2
    \right].
\]
Следовательно,
\[
    (\Delta x)^2
    =
    \langle\hat x^2\rangle
    -
    \langle\hat x\rangle^2
    =
    \boxed{
        \frac{\hbar}{2m\omega}
    }.
\]
То есть
\[
    \boxed{
        \Delta x
        =
        \sqrt{\frac{\hbar}{2m\omega}}
    }.
\]

Аналогично для импульса
\[
    \boxed{
    \langle\hat p\rangle
    =
    -i
    \sqrt{\frac{m\hbar\omega}{2}}
    \left(
        \alpha-\alpha^*
    \right)
    }.
\]
Для квадрата импульса
\begin{align*}
    \langle\hat p^2\rangle
    &=
    -\frac{m\hbar\omega}{2}
    \left\langle
        \hat a^2
        -
        \hat a\hat a^\dagger
        -
        \hat a^\dagger\hat a
        +
        (\hat a^\dagger)^2
    \right\rangle
    \\
    &=
    \frac{m\hbar\omega}{2}
    \left[
        -\alpha^2
        -
        (\alpha^*)^2
        +
        2|\alpha|^2
        +
        1
    \right].
\end{align*}
При этом
\[
    \langle\hat p\rangle^2
    =
    \frac{m\hbar\omega}{2}
    \left[
        -\alpha^2
        -
        (\alpha^*)^2
        +
        2|\alpha|^2
    \right].
\]
Следовательно,
\[
    (\Delta p)^2
    =
    \langle\hat p^2\rangle
    -
    \langle\hat p\rangle^2
    =
    \boxed{
        \frac{m\hbar\omega}{2}
    },
\]
то есть
\[
    \boxed{
        \Delta p
        =
        \sqrt{\frac{m\hbar\omega}{2}}
    }.
\]

Перемножая неопределённости, получаем
\[
    \boxed{
    \Delta x\,\Delta p
    =
    \sqrt{\frac{\hbar}{2m\omega}}
    \sqrt{\frac{m\hbar\omega}{2}}
    =
    \frac{\hbar}{2}
    }.
\]

С другой стороны,
\[
    [\hat x,\hat p]=i\hbar\hat I,
\]
поэтому общее соотношение неопределённостей даёт
\[
    \Delta x\,\Delta p
    \geq
    \frac12
    \left|
        \langle[\hat x,\hat p]\rangle
    \right|
    =
    \frac{\hbar}{2}.
\]
Следовательно, для когерентного состояния неравенство обращается в
равенство:
\[
    \boxed{
        \Delta x\,\Delta p
        =
        \frac{\hbar}{2}
    }.
\]

Таким образом, когерентные состояния являются состояниями с минимальной
неопределённостью координаты и импульса. При этом величины
$\Delta x$ и $\Delta p$ не зависят от параметра $\alpha$; изменение
$\alpha$ изменяет средние значения координаты и импульса, но не ширину
их квантовых распределений.

% =========================================================
\section{Лекции 17--18}
% =========================================================

\subsection{Эволюция в квантовой механике. Картины Шрёдингера и Гейзенберга}

Пусть состояние квантовой системы в момент времени $t_0$ задаётся
вектором $|\psi(t_0)\rangle$. Его эволюцию к моменту времени $t$
опишем оператором эволюции:
\[
    \boxed{
        |\psi(t)\rangle
        =
        \hat U(t,t_0)|\psi(t_0)\rangle
    }.
\]

Скалярное произведение состояний при эволюции должно сохраняться:
\[
    \langle\psi(t)|\psi(t)\rangle
    =
    \langle\psi(t_0)|
        \hat U^\dagger(t,t_0)\hat U(t,t_0)
    |\psi(t_0)\rangle.
\]
Поэтому
\[
    \boxed{
        \hat U^\dagger(t,t_0)\hat U(t,t_0)=\hat I
    }.
\]
Кроме того,
\[
    \hat U(t,t_0)\hat U^\dagger(t,t_0)=\hat I.
\]
Следовательно, оператор эволюции является унитарным:
\[
    \boxed{
        \hat U^{-1}(t,t_0)=\hat U^\dagger(t,t_0)
    }.
\]

Унитарные операторы эволюции обладают групповым свойством. Если
$t_0<t_1<t_2$, то
\[
    \boxed{
        \hat U(t_2,t_1)\hat U(t_1,t_0)
        =
        \hat U(t_2,t_0)
    },
\]
а при отсутствии временного промежутка
\[
    \boxed{
        \hat U(t_0,t_0)=\hat I
    }.
\]

Таким образом, динамика замкнутой квантовой системы описывается
унитарным оператором эволюции $\hat U(t,t_0)$.


% =========================================================
\subsubsection{Инфинитезимальный оператор эволюции}
% =========================================================

Рассмотрим эволюцию на малом промежутке времени $dt$:
\[
    \hat U(t+dt,t)
    =
    \hat I+\hat K(t)\,dt+\ldots
\]
Из унитарности
\[
    \hat U^\dagger(t+dt,t)\hat U(t+dt,t)=\hat I
\]
следует
\[
    \bigl(\hat I+\hat K^\dagger dt+\ldots\bigr)
    \bigl(\hat I+\hat Kdt+\ldots\bigr)
    =
    \hat I.
\]
В первом порядке по $dt$
\[
    \hat K+\hat K^\dagger=0.
\]
Следовательно, $\hat K$ является антиэрмитовым оператором. Его удобно
представить в виде
\[
    \boxed{
        \hat K(t)
        =
        -\frac{i}{\hbar}\hat H(t)
    },
\]
где
\[
    \boxed{
        \hat H^\dagger(t)=\hat H(t)
    }.
\]
Самосопряжённый оператор $\hat H$ является гамильтонианом системы.

Из композиционного свойства оператора эволюции
\[
    \hat U(t+dt,t_0)
    =
    \hat U(t+dt,t)\hat U(t,t_0)
\]
получаем
\[
    \hat U(t+dt,t_0)
    =
    \left(
        \hat I-\frac{i}{\hbar}\hat H(t)dt
    \right)
    \hat U(t,t_0)+o(dt).
\]
Следовательно,
\[
    \boxed{
        \frac{\partial}{\partial t}\hat U(t,t_0)
        =
        -\frac{i}{\hbar}
        \hat H(t)\hat U(t,t_0)
    },
\]
или
\[
    \boxed{
        i\hbar
        \frac{\partial}{\partial t}\hat U(t,t_0)
        =
        \hat H(t)\hat U(t,t_0)
    }.
\]

Применяя это равенство к $|\psi(t_0)\rangle$, получаем
\[
    \boxed{
        i\hbar
        \frac{\partial}{\partial t}
        |\psi(t)\rangle
        =
        \hat H(t)|\psi(t)\rangle
    }.
\]
Это уравнение Шрёдингера в абстрактной форме.


% =========================================================
\subsubsection{Координатное представление. Пропагатор}
% =========================================================

В координатном представлении
\[
    \psi(\vec r,t)
    =
    \langle\vec r|\psi(t)\rangle.
\]
Вставляя разложение единицы между оператором эволюции и начальным
состоянием, получаем
\[
\begin{aligned}
    \psi(\vec r,t)
    &=
    \langle\vec r|
        \hat U(t,t_0)
    |\psi(t_0)\rangle
    \\
    &=
    \int d^3r'\,
        \langle\vec r|
            \hat U(t,t_0)
        |\vec r'\rangle
        \langle\vec r'|\psi(t_0)\rangle.
\end{aligned}
\]
Обозначим
\[
    \boxed{
        G(\vec r,t;\vec r',t_0)
        =
        \langle\vec r|
            \hat U(t,t_0)
        |\vec r'\rangle
    }.
\]
Функция $G$ называется пропагатором. Тогда
\[
    \boxed{
    \psi(\vec r,t)
    =
    \int d^3r'\,
        G(\vec r,t;\vec r',t_0)
        \psi(\vec r',t_0)
    }.
\]


% =========================================================
\subsubsection{Стационарный гамильтониан}
% =========================================================

Если гамильтониан не зависит от времени,
\[
    \hat H(t)=\hat H,
\]
то решение уравнения для оператора эволюции имеет вид
\[
    \boxed{
        \hat U(t,t_0)
        =
        \exp\left[
            -\frac{i}{\hbar}
            \hat H(t-t_0)
        \right]
    }.
\]

Пусть спектральное разложение гамильтониана записано как
\[
    \hat I=\sum_E\hat P_E,
    \qquad
    \hat H=\sum_EE\hat P_E.
\]
Тогда
\[
    |\psi\rangle
    =
    \sum_E\hat P_E|\psi\rangle
\]
и
\[
    \boxed{
        \hat U(t,t_0)
        =
        \sum_E
        e^{-\frac{i}{\hbar}E(t-t_0)}
        \hat P_E
    }.
\]
Следовательно,
\[
    \boxed{
    |\psi(t)\rangle
    =
    \sum_E
        e^{-\frac{i}{\hbar}E(t-t_0)}
        \hat P_E|\psi(t_0)\rangle
    }.
\]
Каждая энергетическая компонента состояния приобретает собственный
фазовый множитель.

Спектральные вероятности энергии при такой эволюции не изменяются:
если
\[
    d\mu_\psi(E)
    =
    \langle\psi|d\hat P_E|\psi\rangle,
\]
то унитарная эволюция, порождённая самим $\hat H$, не меняет
распределение энергии.

При наличии вырождения можно записать
\[
    \hat P_E
    =
    \sum_{a_n}
        |E_n,a_n\rangle
        \langle E_n,a_n|,
\]
и поэтому
\[
    |\psi\rangle
    =
    \sum_{E_n,a_n}
        |E_n,a_n\rangle
        \langle E_n,a_n|\psi\rangle.
\]
Вероятность обнаружить систему в состоянии $|E_n,a_n\rangle$ равна
\[
    \boxed{
        |\langle E_n,a_n|\psi\rangle|^2
    }.
\]


% =========================================================
\subsubsection{Пример: эволюция когерентного состояния}
% =========================================================

Для гармонического осциллятора когерентное состояние имеет вид
\[
    |\alpha\rangle
    =
    C
    \sum_{n=0}^{\infty}
        \frac{\alpha^n}{\sqrt{n!}}|n\rangle,
    \qquad
    \hat a|\alpha\rangle
    =
    \alpha|\alpha\rangle,
    \qquad
    \alpha\in\mathbb C.
\]
Гамильтониан равен
\[
    \hat H
    =
    \hbar\omega
    \left(
        \hat n+\frac12
    \right).
\]
Пусть в начальный момент
\[
    |\psi(0)\rangle=|\alpha\rangle.
\]
Тогда
\begin{align*}
    |\psi(t)\rangle
    &=
    C\sum_{n=0}^{\infty}
        \frac{\alpha^n}{\sqrt{n!}}
        e^{-i\omega(n+\frac12)t}
        |n\rangle
    \\
    &=
    e^{-i\omega t/2}
    C\sum_{n=0}^{\infty}
        \frac{
            \left(
                \alpha e^{-i\omega t}
            \right)^n
        }{\sqrt{n!}}
        |n\rangle.
\end{align*}
Следовательно,
\[
    \boxed{
        |\psi(t)\rangle
        =
        e^{-i\omega t/2}
        |\alpha e^{-i\omega t}\rangle
    }.
\]
Общий фазовый множитель $e^{-i\omega t/2}$ не влияет на вероятности,
поэтому когерентное состояние при эволюции остаётся когерентным, а его
параметр вращается:
\[
    \alpha(t)=\alpha e^{-i\omega t}.
\]

Операторы координаты и импульса выражаются через
$\hat a,\hat a^\dagger$ как
\[
    \hat x
    =
    \sqrt{\frac{\hbar}{m\omega}}\,
    \frac{\hat a+\hat a^\dagger}{\sqrt2},
\]
\[
    \hat p
    =
    \sqrt{m\hbar\omega}\,
    \frac{\hat a-\hat a^\dagger}{i\sqrt2}.
\]
Поэтому
\[
    \boxed{
    \overline{x}(t)
    =
    \langle\psi(t)|\hat x|\psi(t)\rangle
    =
    \sqrt{\frac{\hbar}{m\omega}}\,
    \frac{
        \alpha e^{-i\omega t}
        +
        \alpha^*e^{i\omega t}
    }{\sqrt2}
    }.
\]
Если
\[
    \alpha=|\alpha|e^{i\varphi_\alpha},
\]
то
\[
    \overline{x}(t)
    =
    \sqrt{\frac{2\hbar}{m\omega}}\,
    |\alpha|
    \cos(\omega t-\varphi_\alpha).
\]

Аналогично,
\[
    \boxed{
    \overline{p}(t)
    =
    \sqrt{m\hbar\omega}\,
    \frac{
        \alpha e^{-i\omega t}
        -
        \alpha^*e^{i\omega t}
    }{i\sqrt2}
    }.
\]
То есть средние координата и импульс совершают гармоническое движение,
как координата и импульс классического осциллятора.


% =========================================================
\subsubsection{Картины Шрёдингера и Гейзенберга}
% =========================================================

В картине Шрёдингера эволюционирует вектор состояния:
\[
    \boxed{
        |\psi(t)\rangle
        =
        \hat U(t,t_0)|\psi(t_0)\rangle
    }.
\]
Для оператора $\hat A$ среднее значение равно
\[
\begin{aligned}
    \langle\hat A\rangle_t
    &=
    \langle\psi(t)|
        \hat A
    |\psi(t)\rangle
    \\
    &=
    \langle\psi(t_0)|
        \hat U^\dagger(t,t_0)
        \hat A
        \hat U(t,t_0)
    |\psi(t_0)\rangle.
\end{aligned}
\]

В картине Гейзенберга состояние фиксируется:
\[
    \boxed{
        |\psi\rangle_{\mathrm H}
        =
        |\psi(t_0)\rangle
    },
\]
а зависимость от времени переносится на оператор:
\[
    \boxed{
        \hat A_{\mathrm H}(t)
        =
        \hat U^\dagger(t,t_0)
        \hat A_{\mathrm S}
        \hat U(t,t_0)
    }.
\]
Обе картины дают одинаковые средние значения физических величин.


% =========================================================
\subsection{Дифференцирование операторов по времени}
% =========================================================

Рассмотрим оператор в картине Гейзенберга:
\[
    \hat A_{\mathrm H}(t)
    =
    \hat U^\dagger(t,t_0)
    \hat A_{\mathrm S}(t)
    \hat U(t,t_0).
\]
Если оператор $\hat A_{\mathrm S}$ не имеет явной зависимости от
времени, то
\[
\begin{aligned}
    \frac{d\hat A_{\mathrm H}}{dt}
    &=
    \frac{i}{\hbar}
    \hat U^\dagger
    \hat H
    \hat A_{\mathrm S}
    \hat U
    -
    \frac{i}{\hbar}
    \hat U^\dagger
    \hat A_{\mathrm S}
    \hat H
    \hat U
    \\
    &=
    \frac{i}{\hbar}
    \left[
        \hat H_{\mathrm H},
        \hat A_{\mathrm H}
    \right].
\end{aligned}
\]
Таким образом,
\[
    \boxed{
        \frac{d\hat A_{\mathrm H}}{dt}
        =
        \frac{i}{\hbar}
        [\hat H_{\mathrm H},\hat A_{\mathrm H}]
    }.
\]

Если же $\hat A_{\mathrm S}$ явно зависит от времени, возникает
дополнительный член:
\[
    \boxed{
    \frac{d\hat A_{\mathrm H}}{dt}
    =
    \left(
        \frac{\partial\hat A}{\partial t}
    \right)_{\mathrm H}
    +
    \frac{i}{\hbar}
    [\hat H_{\mathrm H},\hat A_{\mathrm H}]
    },
\]
где
\[
    \boxed{
    \left(
        \frac{\partial\hat A}{\partial t}
    \right)_{\mathrm H}
    =
    \hat U^\dagger
    \left(
        \frac{\partial\hat A_{\mathrm S}}{\partial t}
    \right)
    \hat U
    }.
\]

Соответственно, для среднего значения
\[
    \boxed{
    \frac{d}{dt}
    \langle\hat A\rangle
    =
    \left\langle
        \frac{\partial\hat A}{\partial t}
    \right\rangle
    +
    \frac{i}{\hbar}
    \langle[\hat H,\hat A]\rangle
    }.
\]


% =========================================================
\subsubsection{Теорема Эренфеста}
% =========================================================

Рассмотрим гамильтониан
\[
    \boxed{
        \hat H
        =
        \frac{\hat{\vec p}^{\,2}}{2m}
        +
        V(\hat{\vec r})
    }.
\]
Для координаты
\[
    \dot{\hat r}_\alpha
    =
    \frac{i}{\hbar}
    [\hat H,\hat r_\alpha].
\]
Поскольку
\[
    [\hat p_\beta,\hat r_\alpha]
    =
    -i\hbar\delta_{\alpha\beta},
\]
получаем
\[
    \boxed{
        \dot{\hat{\vec r}}
        =
        \frac{\hat{\vec p}}{m}
    }.
\]

Для импульса
\[
    \boxed{
        \dot{\hat{\vec p}}
        =
        -\nabla V(\hat{\vec r})
    }.
\]
Следовательно,
\[
    \boxed{
        \frac{d}{dt}
        \langle\hat{\vec r}\rangle
        =
        \frac{\langle\hat{\vec p}\rangle}{m}
    },
\]
\[
    \boxed{
        \frac{d}{dt}
        \langle\hat{\vec p}\rangle
        =
        -\langle\nabla V\rangle
    }.
\]
Эти равенства составляют содержание теоремы Эренфеста. Они похожи на
классические уравнения движения, однако вообще говоря
\[
    \langle\nabla V(\hat{\vec r})\rangle
    \neq
    \nabla V(\langle\hat{\vec r}\rangle).
\]


% =========================================================
\subsubsection{Пример 1: свободное движение}
% =========================================================

Для свободной частицы в одном измерении
\[
    \hat H=\frac{\hat p^2}{2m}.
\]
Уравнения Гейзенберга дают
\[
    \frac{d\hat p}{dt}=0,
    \qquad
    \frac{d\hat x}{dt}
    =
    \frac{\hat p}{m}.
\]
Поэтому
\[
    \boxed{
        \hat p(t)=\hat p(0)=\mathrm{const}
    },
\]
\[
    \boxed{
        \hat x(t)
        =
        \hat x(0)
        +
        \frac{\hat p}{m}t
    }.
\]

Рассмотрим коммутатор координат в различные моменты времени:
\[
\begin{aligned}
    [\hat x(t),\hat x(0)]
    &=
    \frac{t}{m}
    [\hat p,\hat x(0)]
    \\
    &=
    -\frac{i\hbar t}{m}\hat I.
\end{aligned}
\]
Следовательно, из соотношения неопределённостей
\[
    \boxed{
        \Delta x(t)\,\Delta x(0)
        \geq
        \frac{\hbar t}{2m}
    }.
\]
Отсюда
\[
    \boxed{
        \Delta x(t)
        \geq
        \frac{\hbar t}
        {2m\,\Delta x(0)}
    }.
\]
Таким образом, локализованный волновой пакет свободной частицы со
временем расплывается.


% =========================================================
\subsubsection{Пример 2: соотношение неопределённостей энергия--время}
% =========================================================

Пусть оператор $\hat A$ не имеет явной зависимости от времени. Тогда
\[
    \frac{d\hat A}{dt}
    =
    \frac{i}{\hbar}[\hat H,\hat A].
\]
Соотношение неопределённостей для $\hat A$ и энергии даёт
\[
    \Delta A\,\Delta E
    \geq
    \frac12
    \left|
        \langle[\hat A,\hat H]\rangle
    \right|.
\]
Поскольку
\[
    [\hat A,\hat H]
    =
    i\hbar\frac{d\hat A}{dt},
\]
получаем
\[
    \boxed{
        \Delta A\,\Delta E
        \geq
        \frac{\hbar}{2}
        \left|
            \frac{d}{dt}
            \langle\hat A\rangle
        \right|
    }.
\]

Введём характерное время изменения наблюдаемой $\hat A$:
\[
    \boxed{
        \Delta\tau
        =
        \frac{\Delta A}
        {
            \left|
                \dfrac{d}{dt}
                \langle\hat A\rangle
            \right|
        }
    }.
\]
Тогда
\[
    \boxed{
        \Delta E\,\Delta\tau
        \geq
        \frac{\hbar}{2}
    }.
\]
Это соотношение неопределённости энергия--время в форме
Мандельштама--Тамма.


% =========================================================
\subsubsection{Пример 3: теорема вириала}
% =========================================================

Рассмотрим связанное стационарное состояние системы с гамильтонианом
\[
    \hat H
    =
    \frac{\hat{\vec p}^{\,2}}{2m}
    +
    V(\hat{\vec r}).
\]
В стационарном состоянии средние значения не зависят от времени, поэтому
\[
    \frac{d}{dt}
    \langle
        \hat{\vec r}\cdot\hat{\vec p}
    \rangle
    =
    0.
\]
С другой стороны,
\[
\begin{aligned}
    \frac{d}{dt}
    (\hat{\vec r}\cdot\hat{\vec p})
    &=
    \dot{\hat{\vec r}}\cdot\hat{\vec p}
    +
    \hat{\vec r}\cdot\dot{\hat{\vec p}}
    \\
    &=
    \frac{\hat{\vec p}^{\,2}}{m}
    -
    \hat{\vec r}\cdot\nabla V.
\end{aligned}
\]
Поскольку
\[
    \frac{\hat{\vec p}^{\,2}}{m}=2\hat T,
\]
получаем
\[
    \boxed{
        2\langle T\rangle
        =
        \left\langle
            \vec r\cdot\nabla V
        \right\rangle
    }.
\]
Это квантово-механическая теорема вириала.

Если потенциал является однородной функцией степени $n$,
\[
    V(\lambda\vec r)=\lambda^nV(\vec r),
\]
то
\[
    \vec r\cdot\nabla V=nV.
\]
Следовательно,
\[
    \boxed{
        2\langle T\rangle
        =
        n\langle V\rangle
    }.
\]

Для гармонического осциллятора
\[
    V\propto r^2,
    \qquad
    n=2,
\]
поэтому
\[
    \boxed{
        \langle T\rangle
        =
        \langle V\rangle
    }.
\]

Для кулоновского потенциала
\[
    V\propto-\frac1r,
    \qquad
    n=-1,
\]
следовательно,
\[
    2\langle T\rangle
    =
    -\langle V\rangle.
\]
При этом
\[
    E=\langle T\rangle+\langle V\rangle,
\]
откуда
\[
    \boxed{
        \langle T\rangle=-E,
        \qquad
        \langle V\rangle=2E
    }.
\]


% =========================================================
\subsection{Симметрии и законы сохранения}
% =========================================================

Рассмотрим преобразование квантовых состояний
\[
    |\psi\rangle
    \longrightarrow
    |\psi'\rangle.
\]
Физическая симметрия должна сохранять вероятности переходов:
\[
    \boxed{
        |\langle\varphi|\psi\rangle|^2
        =
        |\langle\varphi'|\psi'\rangle|^2
    }.
\]

Согласно теореме Вигнера такое преобразование реализуется либо
унитарным линейным, либо антиунитарным антилинейным оператором.

Для антилинейного оператора
\[
    \boxed{
    \hat U
    \left(
        \alpha|\psi_1\rangle
        +
        \beta|\psi_2\rangle
    \right)
    =
    \alpha^*\hat U|\psi_1\rangle
    +
    \beta^*\hat U|\psi_2\rangle
    }.
\]
Простейший пример антилинейного преобразования в координатном
представлении --- комплексное сопряжение:
\[
    \psi(\xi)\longrightarrow\psi^*(\xi).
\]

Ниже основное внимание будем уделять унитарным преобразованиям.


% =========================================================
\subsubsection{Преобразование наблюдаемых}
% =========================================================

Пусть
\[
    |\psi'\rangle=\hat U|\psi\rangle.
\]
Среднее значение наблюдаемой должно оставаться тем же физическим
числом:
\[
    \langle\psi|\hat A|\psi\rangle
    =
    \langle\psi'|\hat A'|\psi'\rangle.
\]
Подставляя $|\psi'\rangle=\hat U|\psi\rangle$, получаем
\[
    \langle\psi|
        \hat U^\dagger\hat A'\hat U
    |\psi\rangle
    =
    \langle\psi|\hat A|\psi\rangle.
\]
Следовательно,
\[
    \boxed{
        \hat A
        =
        \hat U^\dagger\hat A'\hat U
    },
\]
или эквивалентно
\[
    \boxed{
        \hat A'
        =
        \hat U\hat A\hat U^\dagger
    }.
\]

Таким образом, можно рассматривать либо активное преобразование
состояния,
\[
    |\psi\rangle\longrightarrow\hat U|\psi\rangle,
\]
либо соответствующее пассивное преобразование базиса. Физические
результаты в обоих описаниях совпадают.


% =========================================================
\subsubsection{Динамические и непрерывные симметрии}
% =========================================================

Если преобразование является симметрией динамики, гамильтониан
инвариантен:
\[
    \boxed{
        \hat U\hat H\hat U^\dagger
        =
        \hat H
    }.
\]
Эквивалентно,
\[
    \boxed{
        [\hat H,\hat U]=0
    }.
\]

Для непрерывной однопараметрической группы преобразований можно
записать
\[
    \boxed{
        \hat U(\alpha)
        =
        \exp(-i\alpha\hat K)
    },
\]
где генератор $\hat K$ является самосопряжённым:
\[
    \hat K^\dagger=\hat K.
\]
Если
\[
    [\hat H,\hat U(\alpha)]=0
\]
для всех $\alpha$, то
\[
    \boxed{
        [\hat H,\hat K]=0
    }.
\]
Следовательно,
\[
    \frac{d\hat K}{dt}=0.
\]
То есть непрерывной симметрии соответствует сохраняющаяся величина.
Это квантово-механическая форма связи между симметриями и законами
сохранения.


% =========================================================
\subsubsection{Пространственные трансляции}
% =========================================================

В одном измерении оператор трансляции на расстояние $a$ определяется
условием
\[
    \boxed{
        \hat U_a|x\rangle
        =
        |x+a\rangle
    }.
\]
Тогда
\begin{align*}
    \hat U_a|\psi\rangle
    &=
    \int dx\,
        |x+a\rangle
        \langle x|\psi\rangle
    \\
    &=
    \int dx\,
        |x\rangle
        \langle x-a|\psi\rangle.
\end{align*}
Следовательно, в координатном представлении
\[
    \boxed{
        \psi'(x)
        =
        \langle x|\hat U_a|\psi\rangle
        =
        \psi(x-a)
    }.
\]

Рассмотрим действие коммутатора $\hat U_a$ и $\hat x$:
\begin{align*}
    [\hat U_a,\hat x]|x\rangle
    &=
    x|x+a\rangle
    -
    (x+a)|x+a\rangle
    \\
    &=
    -a|x+a\rangle.
\end{align*}
Поэтому
\[
    \boxed{
        [\hat x,\hat U_a]
        =
        a\hat U_a
    }.
\]

Трансляции обладают групповым свойством
\[
    \boxed{
        \hat U_a\hat U_b
        =
        \hat U_{a+b}
    }.
\]
Оператор трансляции можно представить через импульс:
\[
    \boxed{
        \hat U_a
        =
        \exp\left(
            -\frac{i}{\hbar}\hat p\,a
        \right)
    }.
\]
Таким образом, импульс является генератором пространственных
трансляций.

Для бесконечно малого $da$
\[
    \hat U_{da}
    =
    \hat I
    -
    \frac{i}{\hbar}\hat p\,da
    +
    \ldots
\]
и
\[
    [\hat x,\hat U_{da}]
    =
    da\,\hat I.
\]
Отсюда снова получаем
\[
    \boxed{
        [\hat x,\hat p]
        =
        i\hbar\hat I
    }.
\]

То же самое можно увидеть непосредственно из преобразования волновой
функции:
\[
    \psi(x-da)
    =
    \psi(x)
    -
    \frac{i}{\hbar}
    da\,
    \langle x|\hat p|\psi\rangle.
\]
С другой стороны,
\[
    \psi(x-da)
    =
    \psi(x)-da\,\frac{\partial\psi}{\partial x}+\ldots
\]
Поэтому
\[
    \boxed{
        \langle x|\hat p|\psi\rangle
        =
        -i\hbar
        \frac{\partial}{\partial x}
        \langle x|\psi\rangle
    }.
\]


% =========================================================
\subsubsection{Трансляции во времени}
% =========================================================

Для стационарного гамильтониана временная эволюция имеет вид
\[
    \boxed{
        \hat U(t,t_0)
        =
        \exp\left[
            -\frac{i}{\hbar}
            \hat H(t-t_0)
        \right]
    }.
\]
Следовательно, гамильтониан является генератором трансляций во времени:
\[
    |\psi(t_0)\rangle
    \longrightarrow
    |\psi(t_0+t)\rangle.
\]


% =========================================================
\subsubsection{Дискретная симметрия: инверсия}
% =========================================================

Рассмотрим в одном измерении оператор пространственной инверсии
$\hat\Pi_x$:
\[
    \boxed{
        \hat\Pi_x|x\rangle
        =
        |-x\rangle
    }.
\]
Двукратное применение инверсии возвращает исходное состояние:
\[
    \boxed{
        \hat\Pi_x^2=\hat I
    }.
\]
Кроме того, оператор инверсии унитарен:
\[
    \hat\Pi_x^\dagger\hat\Pi_x=\hat I.
\]
Так как
\[
    \hat\Pi_x^{-1}=\hat\Pi_x,
\]
получаем
\[
    \boxed{
        \hat\Pi_x^\dagger
        =
        \hat\Pi_x
    }.
\]
То есть оператор пространственной инверсии одновременно унитарен и
самосопряжён.

Его собственные значения удовлетворяют
\[
    \pi_x^2=1,
\]
поэтому
\[
    \boxed{
        \pi_x=\pm1
    }.
\]
Состояния с
\[
    \hat\Pi_x|e\rangle=|e\rangle
\]
называются чётными, а состояния с
\[
    \hat\Pi_x|o\rangle=-|o\rangle
\]
--- нечётными.

В трёхмерном пространстве полная пространственная инверсия действует
как
\[
    \vec r\longrightarrow-\vec r
\]
и может быть записана через инверсии отдельных координат:
\[
    \boxed{
        \hat\Pi
        =
        \hat\Pi_x\hat\Pi_y\hat\Pi_z
    }.
\]

Если в одномерной задаче
\[
    \boxed{
        [\hat H,\hat\Pi_x]=0
    },
\]
то чётность сохраняется во времени. Например, это имеет место для
гамильтониана с чётным потенциалом
\[
    V(x)=V(-x).
\]


% =========================================================
\subsubsection{Пример: инверсионные состояния молекулы \texorpdfstring{$NH_3$}{NH3}}
% =========================================================

В молекуле аммиака $NH_3$ существуют два эквивалентных положения
атома азота относительно плоскости атомов водорода. Обозначим
соответствующие локализованные состояния как
\[
    |n\rangle,
    \qquad
    |v\rangle.
\]
Отдельно они не являются собственными состояниями полного
гамильтониана:
\[
    \hat H|n\rangle\neq E|n\rangle,
    \qquad
    \hat H|v\rangle\neq E|v\rangle.
\]

В базисе $\{|n\rangle,|v\rangle\}$ гамильтониан можно представить в
форме
\[
    \boxed{
        \hat H
        =
        \begin{pmatrix}
            E & V\\
            V & E
        \end{pmatrix}
    }.
\]
Его собственными состояниями являются симметричная и
антисимметричная комбинации
\[
    \boxed{
        |+\rangle
        =
        \frac{1}{\sqrt2}
        \left(
            |n\rangle+|v\rangle
        \right)
    },
\]
\[
    \boxed{
        |-\rangle
        =
        \frac{1}{\sqrt2}
        \left(
            |n\rangle-|v\rangle
        \right)
    }.
\]
Соответствующие энергии равны
\[
    \boxed{
        E_+=E+V,
        \qquad
        E_-=E-V
    }.
\]

Таким образом, туннелирование между двумя эквивалентными
конфигурациями снимает вырождение и приводит к расщеплению
энергетического уровня. Инверсионный переход между этими уровнями в
молекуле аммиака был использован в первом мазере.

\end{document}
